\documentclass[conference]{IEEEtran}
\IEEEoverridecommandlockouts
% The preceding line is only needed to identify funding in the first footnote. If that is unneeded, please comment it out.
%Template version as of 6/27/2024
\usepackage[slovak,english]{babel}
\usepackage{tikz}
\usetikzlibrary{positioning,arrows.meta}
\usepackage{graphicx}
\usepackage{cite}
\usepackage{amsmath,amssymb,amsfonts}
\usepackage{algorithmic}
\usepackage{float}
\usepackage{booktabs}
\usepackage{textcomp}
\usepackage[T1]{fontenc}
\usepackage[utf8]{inputenc}
\usepackage{stfloats}
\usepackage{lmodern}

\usepackage{xcolor}
\def\BibTeX{{\rm B\kern-.05em{\sc i\kern-.025em b}\kern-.08em
    T\kern-.1667em\lower.7ex\hbox{E}\kern-.125emX}}
\begin{document}
\selectlanguage{slovak}
\title{Bezpečnosť elektronickej pošty – analýza hrozieb a návrh nástroja na detekciu podozrivých e-mailov\\
}
\author{\IEEEauthorblockN{Alexander Ovšonka}
\textit{STU FIIT}\\
xovsonka@stuba.sk}


\maketitle

\begin{abstract}
Táto práca sa zaoberá problematikou bezpečnosti e-mailovej komunikácie a analýzou hrozieb, ktoré sa v nej vyskytujú. V práci sú predstavené základné princípy fungovania e-mailovej komunikácie, formát e-mailových správ a najčastejšie typy útokov, ako sú phishing, spam alebo spoofing.

Súčasťou práce je návrh vlastného riešenia na detekciu podozrivých e-mailov, ktoré kombinuje heuristický analyzátor s prístupmi strojového učenia. Navrhnutý systém využíva extrakciu príznakov z obsahu aj hlavičiek správ a umožňuje klasifikáciu e-mailov do viacerých kategórií.

Navrhnuté riešenie je experimentálne overené na dátach získaných z verejných datasetov aj simulovaných kampaní, pričom výsledky poukazujú na jeho schopnosť efektívne detegovať škodlivé e-mailové správy.
\end{abstract}

\begin{IEEEkeywords}
email security, phishing detection, spam detection, email communication, machine learning, heuristic analysis, cybersecurity, email threats
\end{IEEEkeywords}

\section{Úvod}

Elektronická pošta patrí medzi najrozšírenejšie formy digitálnej komunikácie, no zároveň predstavuje jeden z najvýznamnejších vektorov kybernetických útokov. Útočníci čoraz častejšie využívajú pokročilé techniky sociálneho inžinierstva a sofistikované formy manipulácie s cieľom oklamať používateľa, získať citlivé údaje alebo distribuovať škodlivý softvér, ako napríklad malware či ransomware.

S rastúcou kvalitou a vierohodnosťou podvodných e-mailov sa ich detekcia stáva čoraz náročnejšou nielen pre automatizované systémy, ale aj pre samotných používateľov. Tento trend vedie k potrebe vývoja robustnejších detekčných mechanizmov, ktoré dokážu identifikovať škodlivé správy aj v podmienkach meniacich sa útočných techník.

Cieľom tejto práce je analyzovať problematiku bezpečnosti e-mailovej komunikácie a navrhnúť nástroj na detekciu podozrivých e-mailov. V teoretickej časti sa zameriavame na princípy fungovania e-mailovej komunikácie, jej architektúru a formát správ, ako aj na najčastejšie typy útokov a indikátory podvodných e-mailov. Zároveň predstavujeme základné mechanizmy ochrany používané v praxi.

Praktická časť práce sa venuje návrhu a implementácii vlastného riešenia, ktoré kombinuje heuristický analyzátor s prístupmi strojového učenia. Súčasťou je návrh klasifikačného modelu, extrakcia relevantných príznakov a realizácia experimentov zameraných na overenie presnosti a robustnosti navrhnutého systému. Okrem toho sú analyzované aj existujúce open-source nástroje a ich využitie v kontexte bezpečnosti e-mailovej komunikácie.
\section{Teoretické základy}
\subsection{Architektúra e-mailovej komunikácie}

Elektronická pošta patrí medzi najstaršie a najrozšírenejšie služby internetu. 
Jej fungovanie je založené na súbore štandardizovaných internetových protokolov 
a formátov správ, ktoré umožňujú vytváranie, prenos a doručovanie e-mailových správ 
medzi jednotlivými účastníkmi komunikácie\cite{rfc5598}.

Základom e-mailovej infraštruktúry je protokol \textbf{SMTP} (Simple Mail Transfer Protocol), 
ktorý zabezpečuje prenos e-mailových správ medzi servermi v internete. 
Formát samotnej e-mailovej správy je definovaný štandardom 
\textbf{Internet Message Format}, zatiaľ čo rozšírenie \textbf{MIME (Multipurpose Internet Mail Extensions)} 
rozširuje formát e-mailovej správy tak, aby bolo možné prenášať 
rôzne typy dát, napríklad binárne prílohy, obrázky alebo HTML obsah\cite{rfc5598}.

Architektúra elektronickej pošty je vysoko distribuovaná. 
Na jej fungovaní sa podieľa viacero komponentov a aktérov, ktorí spolupracujú 
pri spracovaní a prenose správy od odosielateľa k príjemcovi. 
Podľa architektúry možno týchto aktérov rozdeliť 
do niekoľkých základných kategórií\cite{rfc5598}.\\


\textbf{Používatelia}

Používateľ predstavuje entitu, ktorá vytvára alebo prijíma e-mailové správy. 
Používateľom môže byť človek, organizácia alebo automatizovaný proces. 
V rámci e-mailovej architektúry možno rozlíšiť viacero rolí používateľov\cite{rfc5598}.

\begin{itemize}
\item \textbf{Author (autor)} – entita zodpovedná za vytvorenie správy a určenie jej príjemcov.
\item \textbf{Recipient (príjemca)} – cieľová entita, ktorej je správa určená.
\item \textbf{Return handler} – špeciálny typ príjemcu, ktorý prijíma systémové notifikácie generované počas prenosu správy (napr. chybové hlásenia o nedoručení).
\item \textbf{Mediator} – komponent alebo služba, ktorá prijíma správu, modifikuje ju alebo redistribuuje medzi ďalších príjemcov. Typickým príkladom sú mailing listy alebo automatické presmerovania.
\end{itemize}

Z pohľadu používateľa sa proces prenosu e-mailu javí ako jednotná služba 
poskytovaná systémom spracovania správ, nazývaným \textbf{Message Handling System (MHS)}. 
V praxi však môže byť táto služba implementovaná viacerými nezávislými systémami 
a organizáciami\cite{rfc5598}.\\

\textbf{Message Handling System}

Predstavuje infraštruktúru, ktorá zabezpečuje 
kompletný prenos správy od autora k príjemcovi. 
Jeho úlohou je vykonať prenos správy spôsobom \textit{store-and-forward}, 
pri ktorom sa správa postupne prenáša medzi jednotlivými servermi, 
kým nie je doručená cieľovému príjemcovi\cite{rfc5598}.

MHS pozostáva z viacerých komponentov\cite{rfc5598}:

\begin{itemize}
\item \textbf{Message User Agent (MUA)} – klientská aplikácia používaná používateľom 
na vytváranie, čítanie a odosielanie e-mailových správ.
\item \textbf{Mail Submission Agent (MSA)} – prijíma správy od MUA a kontroluje ich 
súlad s pravidlami domény a internetovými štandardmi.
\item \textbf{Message Transfer Agent (MTA)} – zodpovedá za prenos správy medzi jednotlivými 
poštovými servermi. Na komunikáciu medzi MTA sa používa protokol SMTP.
\item \textbf{Mail Delivery Agent (MDA)} – zabezpečuje konečné doručenie správy 
do mailboxu príjemcu.
\item \textbf{Message Store} – úložisko správ, z ktorého si používateľ môže 
svoje e-maily načítať prostredníctvom klienta.
\end{itemize}

Prenos správy medzi jednotlivými servermi prebieha postupne, pričom každý server 
správu dočasne uloží a následne ju prepošle ďalšiemu serveru. 
Tento model komunikácie sa označuje ako \textbf{store-and-forward}\cite{rfc5598}.\\

\textbf{Relay a Gateway}

V rámci infraštruktúry MHS môžu vystupovať aj ďalšie komponenty\cite{rfc5598}.

\begin{itemize}
\item \textbf{Relay} – server, ktorý prijíma e-mailovú správu a preposiela ju 
ďalšiemu serveru smerom k cieľovému príjemcovi.
\item \textbf{Gateway} – komponent prepájajúci rôzne e-mailové systémy alebo 
komunikačné protokoly. Gateway môže upravovať obsah správy alebo jej formát 
tak, aby bola kompatibilná s cieľovým systémom.
\end{itemize}

\textbf{Adresovanie e-mailov}

Každý používateľ e-mailovej služby je identifikovaný pomocou e-mailovej adresy. 
Tá má všeobecný tvar\cite{rfc5598}:

\[
local\text{-}part@domain
\]

Adresa sa skladá z dvoch častí\cite{rfc5598}:

\begin{itemize}
\item \textbf{local-part} – identifikátor používateľa v rámci konkrétnej domény,
\item \textbf{domain} – názov domény, ktorá identifikuje poštový server zodpovedný 
za doručenie správy.
\end{itemize}

Doménové meno je globálne jedinečný identifikátor internetového zdroja, 
ktorý sa pomocou systému DNS mapuje na IP adresu servera\cite{rfc5598}.

\textbf{Identifikácia správ}

Každá e-mailová správa môže byť identifikovaná pomocou jedinečných identifikátorov\cite{rfc5598}.

\begin{itemize}
\item \textbf{Message-ID} – jedinečný identifikátor správy používaný na sledovanie 
a referencovanie správy v rámci e-mailovej komunikácie\cite{rfc5598}.
\item \textbf{ENVID} – identifikátor používaný pri sledovaní prenosu správy 
počas jej doručovania v rámci SMTP komunikácie\cite{rfc5598}.
\end{itemize}

Tieto identifikátory umožňujú sledovať stav správy počas jej prenosu 
a uľahčujú diagnostiku problémov pri doručovaní.

\textbf{Asynchrónna komunikácia}

Prenos e-mailových správ prebieha asynchrónne. 
Správa sa postupne presúva medzi jednotlivými servermi, 
ktoré ju ukladajú a následne preposielajú ďalšiemu serveru. 
Tento model komunikácie umožňuje spoľahlivý prenos správ aj v prípade 
dočasnej nedostupnosti niektorých serverov alebo sieťových spojení\cite{rfc5321}.
\textbf{DNS a MX záznamy}

Pri doručovaní e-mailovej správy je potrebné určiť server, ktorý je 
zodpovedný za prijatie správy pre cieľovú doménu. Tento proces je 
realizovaný pomocou systému doménových mien (DNS). Doménové meno 
predstavuje globálny identifikátor internetového zdroja, ktorý je možné 
prostredníctvom DNS preložiť na IP adresu konkrétneho servera\cite{rfc5321}.

Pri doručovaní e-mailov sa využívajú najmä \textbf{MX (Mail Exchange) záznamy}, 
ktoré určujú poštové servery zodpovedné za prijímanie správ pre danú doménu. 
Ak server potrebuje doručiť e-mail na adresu napríklad \textit{user@example.com}, 
najprv z adresy extrahuje doménu \textit{example.com}. Následne vykoná DNS 
dotaz na MX záznamy tejto domény. Výsledkom môže byť jeden alebo viacero 
poštových serverov, pričom každý má priradenú prioritu určujúcu poradie 
použitia\cite{rfc5321}.

Server sa následne pokúsi nadviazať SMTP spojenie s poštovým serverom 
s najvyššou prioritou. Ak sa spojenie nepodarí nadviazať, môže skúsiť 
ďalší server v poradí. V prípade, že doména nemá definovaný MX záznam, 
SMTP špecifikácia umožňuje použiť ako cieľový server host definovaný 
pomocou \textbf{A alebo AAAA DNS záznamu}. Tento mechanizmus zabezpečuje 
spoľahlivé smerovanie e-mailovej komunikácie v distribuovanom prostredí 
internetu\cite{rfc5321}.

\textbf{SMTP komunikácia medzi servermi}

Prenos e-mailových správ medzi poštovými servermi je realizovaný pomocou 
protokolu \textbf{SMTP}. SMTP používa 
model komunikácie klient–server, kde jeden server vystupuje ako klient 
iniciujúci prenos správy a druhý server ako server prijímajúci správu\cite{rfc5321}.

Po nadviazaní TCP spojenia medzi servermi prebieha komunikácia pomocou 
sekvencie textových príkazov a odpovedí. Komunikácia začína identifikáciou 
servera pomocou príkazu \textbf{HELO} alebo jeho rozšírenej verzie 
\textbf{EHLO}. Následne odosielajúci server oznámi adresu odosielateľa 
pomocou príkazu \textbf{MAIL FROM}. Pre každého príjemcu je následne 
použitý príkaz \textbf{RCPT TO}, ktorý identifikuje cieľového adresáta 
správy\cite{rfc5321}.

Po určení všetkých príjemcov je prenesený samotný obsah správy pomocou 
príkazu \textbf{DATA}. Správa obsahuje hlavičky a telo správy vo formáte 
IMF. Po úspešnom prenose správy server potvrdí jej 
prijatie a komunikácia môže byť ukončená príkazom \textbf{QUIT}\cite{rfc5321}.

SMTP využíva model \textit{store-and-forward}, pri ktorom jednotlivé 
poštové servery dočasne ukladajú prijaté správy a následne ich 
preposielajú ďalšiemu serveru smerom k cieľovej doméne. Tento prístup 
umožňuje spoľahlivý prenos správ aj v prípade dočasnej nedostupnosti 
niektorých serverov alebo sieťových spojení\cite{rfc5321}.

Počas prenosu správy pridáva každý poštový server do hlavičky 
správy záznam \textit{Received}, ktorý obsahuje informácie o trase 
doručenia správy. Tieto záznamy umožňujú sledovať cestu e-mailu 
naprieč jednotlivými servermi a sú dôležité napríklad pri analýze 
pôvodu podozrivých alebo phishingových správ\cite{rfc5321}.

\subsection{Formát e-mailovej správy}

Formát e-mailovej správy je definovaný štandardom 
IMF. Tento štandard definuje syntaktickú štruktúru e-mailových správ, 
vrátane ich hlavičiek a tela správy \cite{rfc5322}.

Na najzákladnejšej úrovni je e-mailová správa reprezentovaná ako sekvencia 
ASCII znakov. Správa je rozdelená do riadkov znakov, pričom každý riadok 
je ukončený dvojicou znakov \textbf{CRLF} (Carriage Return a Line Feed). 
Podľa špecifikácie nesmie jeden riadok správy presiahnuť 998 znakov a 
odporúčaná maximálna dĺžka riadku je 78 znakov bez započítania CRLF 
\cite{rfc5322}.

E-mailová správa pozostáva z dvoch hlavných častí:

\begin{itemize}
\item \textbf{Hlavička správy (Header)}
\item \textbf{Telo správy (Body)}
\end{itemize}

Tieto dve časti sú od seba oddelené prázdnym riadkom.

\subsubsection{Hlavička e-mailovej správy}

Hlavička správy obsahuje metadáta popisujúce správu a jej prenos. 
Je tvorená množinou hlavičkových polí, pričom každé pole má tvar:

\[
Field-Name: Field-Body
\]

Názov poľa je tvorený ASCII znakmi a je ukončený dvojbodkou. 
Obsah poľa môže obsahovať text alebo štruktúrované údaje podľa 
syntaktických pravidiel \cite{rfc5322}.

Dôležité hlavičkové polia zahŕňajú najmä \cite{rfc5322}:

\begin{itemize}
\item \textbf{Date} – dátum a čas vytvorenia správy.
\item \textbf{From} – adresa autora správy.
\item \textbf{Sender} – identifikuje odosielateľa v prípade, že správa 
bola odoslaná v mene viacerých autorov.
\item \textbf{Reply-To} – adresa, na ktorú majú byť zasielané odpovede.
\item \textbf{To} – adresa hlavného príjemcu správy.
\item \textbf{Cc} – adresy ďalších príjemcov správy.
\item \textbf{Bcc} – skrytí príjemcovia správy.
\item \textbf{Message-ID} – jedinečný identifikátor správy.
\item \textbf{In-Reply-To} – identifikátor správy, na ktorú sa odpovedá.
\item \textbf{References} – identifikátory súvisiacich správ v konverzácii.
\item \textbf{Subject} – stručný opis obsahu správy.
\end{itemize}

Hlavičkové polia môžu obsahovať štruktúrované aj neštruktúrované údaje. 
V niektorých prípadoch môže byť dlhý obsah hlavičky rozdelený do viacerých 
riadkov pomocou mechanizmu nazývaného \textit{header folding}, ktorý umožňuje 
zalomiť riadok a pokračovať na ďalšom riadku začínajúcom medzerou alebo 
tabulátorom \cite{rfc5322}.

\subsubsection{Telo správy}

Telo správy obsahuje samotný obsah e-mailu určený pre príjemcu. 
Je tvorené sekvenciou 
textových riadkov kódovaných v ASCII. Každý riadok musí byť ukončený 
sekvenciou CRLF a nesmie prekročiť maximálnu dĺžku 998 znakov \cite{rfc5322}.

Moderné e-mailové správy však často obsahujú aj iné typy dát, napríklad 
formátovaný text, obrázky alebo binárne prílohy. Na podporu týchto 
funkcionalít bol definovaný štandard \textbf{MIME} \cite{rfc2045}.

\subsubsection{MIME rozšírenia}

MIME rozširuje základný formát e-mailovej správy o mechanizmy, ktoré 
umožňujú prenášať rôzne typy dát a viacero častí správy v jednom e-maile. 
MIME zavádza niekoľko nových hlavičkových polí, medzi ktoré patria najmä\cite{rfc2045}:

\begin{itemize}
\item \textbf{MIME-Version} – identifikuje správu ako MIME kompatibilnú.
\item \textbf{Content-Type} – určuje typ obsahu správy alebo jej časti 
(napr. text/plain, text/html alebo image/jpeg).
\item \textbf{Content-Transfer-Encoding} – určuje spôsob kódovania 
obsahu správy, napríklad base64 alebo quoted-printable.
\end{itemize}

Hlavička \textit{Content-Type} umožňuje špecifikovať typ prenášaných 
dát pomocou tzv. mediálnych typov. Tie sú definované vo formáte 
\textit{type/subtype}\cite{rfc2045}. Napríklad:

\begin{itemize}
\item text/plain – jednoduchý text,
\item text/html – HTML obsah správy,
\item image/jpeg – obrázok vo formáte JPEG.
\end{itemize}

MIME tiež umožňuje vytvárať viacdielne správy pomocou typu 
\textbf{multipart}. Takáto správa môže obsahovať viacero častí, 
napríklad textovú verziu správy a HTML verziu alebo rôzne prílohy.\cite{rfc2045}

\subsubsection{Kódovanie dát}

Pri prenose e-mailov cez SMTP je potrebné zabezpečiť kompatibilitu 
s prenosovými mechanizmami, ktoré pôvodne podporovali len 7-bitové 
ASCII znaky. MIME preto definuje rôzne metódy kódovania dát.\cite{rfc2045}

Medzi najčastejšie používané patria\cite{rfc2045}:

\begin{itemize}
\item \textbf{7bit} – základný ASCII formát bez rozšírených znakov.
\item \textbf{8bit} – umožňuje používať znaky mimo základnej ASCII tabuľky.
\item \textbf{base64} – kódovanie používané najmä pre binárne prílohy.
\item \textbf{quoted-printable} – kódovanie používané pre text s 
rozšírenými znakmi.
\end{itemize}

Tieto mechanizmy umožňujú bezpečný prenos rôznych typov dát 
prostredníctvom e-mailovej infraštruktúry.

\subsection{Bezpečnostné hrozby e-mailovej komunikácie}

Elektronická pošta patrí medzi najčastejšie využívané komunikačné služby internetu, 
avšak zároveň predstavuje jeden z najvýznamnejších vektorov kybernetických útokov. 
Útočníci často zneužívajú e-mailovú komunikáciu na šírenie nevyžiadaných správ, 
phishingových kampaní, podvodných odkazov alebo škodlivého softvéru. 
Tieto hrozby môžu viesť k finančným stratám, poškodeniu reputácie organizácií 
alebo kompromitácii citlivých údajov používateľov \cite{karim2019spam}.

\textbf{Spam}

Spam predstavuje rozosielanie nevyžiadaných správ, ktoré sú vo väčšine prípadov 
distribuované hromadne prostredníctvom e-mailových služieb. Takéto správy sú 
často využívané na marketingové účely, avšak môžu byť zároveň súčasťou 
sofistikovanejších útokov, ktorých cieľom je finančný zisk alebo poškodenie 
používateľov a organizácií \cite{karim2019spam}.

Spamové e-maily môžu obsahovať škodlivé odkazy, podvodné prílohy alebo sa môžu 
pokúšať presvedčiť používateľa, aby vykonal určitú akciu, napríklad poskytol 
citlivé informácie alebo navštívil podvodnú webovú stránku. V praxi je často 
ťažké jednoznačne kategorizovať všetky typy spamových správ, pretože spam 
predstavuje širokú skupinu rôznych typov útokov a podvodných aktivít 
\cite{karim2019spam}.

Motiváciou pre šírenie spamových správ je najčastejšie finančný zisk. 
Útočníci využívajú spam napríklad na šírenie phishingových kampaní, 
distribúciu škodlivého softvéru alebo propagáciu podvodných služieb. 
Podľa niektorých odhadov môžu spameri dosahovať značné finančné zisky 
prostredníctvom masového rozosielania spamových správ \cite{karim2019spam}.

Spamové útoky môžu mať rôzne formy. Medzi najčastejšie patria správy 
obsahujúce škodlivé odkazy, prílohy so škodlivým softvérom alebo pokusy 
o phishing a spoofing. Niektorí útočníci sa dokonca snažia obchádzať 
spamové filtre tým, že škodlivý text ukrývajú v obrázkoch alebo iných 
multimediálnych prvkoch \cite{karim2019spam}.

\textbf{Phishing}

Typický priebeh útoku realizovaného prostredníctvom e-mailovej komunikácie
je znázornený na obrázku \ref{fig:email_attack_model}.
\begin{figure}[h]
\centering
\resizebox{\columnwidth}{!}{%
\begin{tikzpicture}[
    node distance=0.5cm,
    every node/.style={font=\footnotesize},
    block/.style={
        draw,
        rounded corners=4pt,
        align=center,
        minimum width=4.2cm,
        minimum height=0.9cm,
        inner sep=4pt
    },
    myarrow/.style={
        -{Latex[length=2mm]},
        thick
    }
]

\node[block] (attacker) {Útočník};
\node[block, below=of attacker] (email) {Podvodná e-mailová správa\\(spam, phishing alebo spoofing)};
\node[block, below=of email] (delivery) {Doručenie správy\\do schránky používateľa};
\node[block, below=of delivery] (interaction) {Interakcia používateľa\\otvorenie prílohy alebo kliknutie na odkaz};
\node[block, below=of interaction] (execution) {Spustenie škodlivého obsahu\\alebo presmerovanie na podvodnú stránku};
\node[block, below=of execution] (consequence) {Následok útoku\\krádež prihlasovacích údajov, infekcia systému\\alebo zašifrovanie súborov};

\draw[myarrow] (attacker) -- (email);
\draw[myarrow] (email) -- (delivery);
\draw[myarrow] (delivery) -- (interaction);
\draw[myarrow] (interaction) -- (execution);
\draw[myarrow] (execution) -- (consequence);

\end{tikzpicture}%
}
\caption{Model priebehu útoku prostredníctvom e-mailovej komunikácie}
\label{fig:email_attack_model}
\end{figure}

Phishing predstavuje jednu z najrozšírenejších foriem e-mailových útokov. 
Pri tomto type útoku sa útočník snaží oklamať používateľa tým, že sa 
vydáva za dôveryhodnú organizáciu alebo známu osobu. Cieľom phishingu 
je získať citlivé údaje používateľa, napríklad prihlasovacie údaje, 
finančné informácie alebo osobné údaje \cite{karim2019spam}.

Phishingové správy často obsahujú odkazy na podvodné webové stránky, 
ktoré napodobňujú vzhľad legitímnych služieb. Používateľ je následne 
vyzvaný na zadanie citlivých údajov, ktoré sú následne odoslané útočníkovi. 
Phishingové útoky môžu byť realizované aj manipuláciou s hlavičkami 
e-mailovej správy, napríklad úpravou polí \textit{From} alebo údajov 
SMTP obálky tak, aby správa pôsobila dôveryhodne \cite{karim2019spam}.

Útočníci môžu taktiež manipulovať s doménou uvedenou v príkaze 
\textit{HELO/EHLO} počas SMTP komunikácie, čím sa snažia vytvoriť dojem, 
že správa pochádza z legitímneho zdroja \cite{karim2019spam}.

\subsubsection*{Spear phishing}

predstavuje pokročilejšiu formu phishingového útoku, 
ktorá je zameraná na konkrétneho jednotlivca alebo organizáciu. 
Na rozdiel od bežného phishingu sú tieto útoky výrazne personalizované 
a často využívajú informácie získané prostredníctvom sociálneho 
inžinierstva \cite{karim2019spam}.

Útočníci sa pri spear phishingu snažia vytvoriť e-mailovú správu, 
ktorá pôsobí ako legitímna komunikácia od dôveryhodnej osoby alebo 
organizácie. Personalizácia správy môže zahŕňať použitie mena 
príjemcu, informácií o jeho pracovnej pozícii alebo detailov o 
organizácii, v ktorej pracuje. Práve táto personalizácia výrazne 
sťažuje detekciu takýchto útokov pomocou tradičných spamových filtrov 
\cite{karim2019spam}.

\subsubsection*{Whaling}

Whaling je špecifická forma spear phishingu, ktorá je zameraná 
na vysoko postavených zamestnancov organizácie, napríklad členov 
manažmentu alebo vedúcich pracovníkov. Útočníci sa v takýchto 
prípadoch snažia napodobniť komunikáciu od dôveryhodných osôb, 
napríklad od generálneho riaditeľa alebo iného člena vedenia 
organizácie \cite{karim2019spam}.

Obsah takýchto e-mailov často vytvára dojem naliehavej situácie, 
napríklad urgentnej finančnej transakcie alebo riešenia 
zákazníckej sťažnosti. Takéto útoky môžu mať veľmi vážne 
finančné dôsledky. V minulosti bolo zaznamenaných viacero 
incidentov, pri ktorých organizácie prišli o milióny eur v 
dôsledku sofistikovaných phishingových kampaní zameraných 
na manažment spoločností \cite{karim2019spam}.

\textbf{Spoofing odosielateľa}

Spoofing odosielateľa predstavuje techniku, pri ktorej útočník
manipuluje s informáciami v e-mailovej správe tak, aby vytvoril
dojem, že správa pochádza od dôveryhodného zdroja. V praxi ide
najčastejšie o úpravu údajov v hlavičke správy alebo v SMTP
obálke, napríklad manipuláciou s poľom \textit{From} alebo
inými identifikačnými údajmi odosielateľa\cite{karim2019spam}.

Pri takomto útoku sa útočník snaží zamaskovať skutočný pôvod
správy tak, aby sa správa javila ako odoslaná legitímnou
organizáciou alebo známou osobou. Spoofing je často využívaný
v kombinácii s phishingovými útokmi, kde cieľom útočníka je
získať citlivé údaje používateľa alebo ho presvedčiť, aby
vykonal určitú akciu, napríklad otvoril škodlivú prílohu alebo
navštívil podvodnú webovú stránku\cite{karim2019spam}.

Útočníci môžu manipulovať aj s údajmi prenášanými počas SMTP
komunikácie, napríklad s doménou uvedenou v príkaze
\textit{HELO/EHLO}. Takýmto spôsobom sa snažia vytvoriť dojem,
že správa bola odoslaná z legitímneho servera alebo známej
domény \cite{karim2019spam}.

Spoofing je umožnený najmä historickým návrhom protokolu SMTP,
ktorý pôvodne neobsahoval mechanizmy na overovanie identity
odosielateľa. Tento nedostatok bol v priebehu času čiastočne
riešený zavedením autentifikačných mechanizmov, ako sú
Sender Policy Framework (SPF) alebo DomainKeys Identified
Mail (DKIM), ktoré umožňujú overovať autenticitu odosielateľa
\cite{karim2019spam}.

\textbf{Distribúcia malvéru prostredníctvom e-mailu}

Podľa študie sú e-mailové prílohy jedným z najčastejších
spôsobov distribúcie ransomvéru a iného malvéru. Útočníci často
maskujú škodlivé súbory ako bežné dokumenty alebo archívy,
napríklad vo formátoch \texttt{ZIP}, \texttt{PDF}, \texttt{DOC} alebo
\texttt{SCR}. Po otvorení takéhoto súboru sa v systéme používateľa
vykoná škodlivý kód, ktorý môže napríklad stiahnuť ďalší malvér,
získať neoprávnený prístup k systému alebo začať šifrovať súbory
používateľa \cite{aurangzeb2017ransomware}

Jedným z najrozšírenejších typov malvéru distribuovaných prostredníctvom
e-mailu je ransomware. Tento typ škodlivého softvéru po infikovaní
zariadenia zablokuje prístup k systému alebo zašifruje súbory
používateľa a následne požaduje zaplatenie výkupného za ich
obnovenie
\cite{aurangzeb2017ransomware}. 

Po úspešnom infikovaní systému môže malvér nadviazať spojenie
so vzdialeným riadiacim serverom (Command and Control server),
ktorý útočníkom umožňuje kontrolovať infikované zariadenie,
získavať citlivé údaje alebo distribuovať ďalší škodlivý kód.
Takéto infikované zariadenia môžu byť následne zapojené do
botnetov, ktoré útočníci využívajú napríklad na šírenie ďalších
spamových kampaní alebo na vykonávanie distribuovaných
útokov typu Denial-of-Service \cite{aurangzeb2017ransomware}. 



\textbf{Indikátory podozrivých e-mailov}

Pri analýze e-mailovej komunikácie je možné identifikovať viacero
indikátorov, ktoré môžu naznačovať, že správa je podozrivá alebo
súčasťou phishingovej či spamovej kampane. Takéto indikátory sa môžu
vyskytovať v hlavičke správy, v odkazoch obsiahnutých v tele správy
alebo v samotnom obsahu e-mailu. Identifikácia týchto znakov predstavuje
dôležitú súčasť analýzy e-mailovej komunikácie a využíva sa aj v
systémoch detekcie phishingu a spamových správ \cite{jakobsson2007phishing,verma2017phishing}.

\textbf{Indikátory v hlavičkách e-mailu.}
Hlavičky e-mailovej správy obsahujú technické informácie o prenose
správy medzi jednotlivými servermi. Analýza týchto údajov môže odhaliť
nezrovnalosti alebo manipuláciu s identitou odosielateľa. Medzi
časté indikátory patrí napríklad nesúlad medzi adresami uvedenými
v poliach \textit{From} a \textit{Return-Path}, podozrivá alebo
nezvyčajná reťaz serverov v záznamoch \textit{Received} alebo
nezhoda medzi deklarovanou doménou odosielateľa a skutočným
serverom, z ktorého bola správa odoslaná \cite{jakobsson2007phishing}.

\textbf{Indikátory v URL odkazoch.}
Phishingové správy často obsahujú odkazy na podvodné webové stránky.
Tieto odkazy môžu využívať rôzne techniky na maskovanie skutočnej
adresy cieľovej stránky. Medzi typické indikátory patrí použitie
domén, ktoré sú vizuálne podobné známym službám (tzv. typosquatting),
použitie IP adresy namiesto doménového mena alebo využívanie služieb
na skracovanie URL adries, ktoré sťažujú identifikáciu skutočnej
cieľovej stránky \cite{verma2017phishing}.

\textbf{Indikátory v obsahu správy.}
Samotný obsah e-mailovej správy môže taktiež obsahovať znaky
naznačujúce potenciálne škodlivý charakter správy. Útočníci často
využívajú naliehavý alebo manipulatívny jazyk, ktorým sa snažia
prinútiť používateľa konať bez dôkladného overenia správy.
Typickými príkladmi sú výzvy na okamžité zadanie prihlasovacích
údajov, potvrdenie účtu alebo vykonanie finančnej transakcie
\cite{jakobsson2007phishing}.

\textbf{Podozrivé prílohy.}
Ďalším častým indikátorom sú prílohy obsahujúce potenciálne
škodlivý kód. Útočníci môžu distribuovať malware prostredníctvom
dokumentov, archívov alebo spustiteľných súborov, ktoré po otvorení
infikujú systém používateľa. Takéto prílohy sú často maskované
ako faktúry, formuláre alebo iné legitímne dokumenty a predstavujú
jednu z najčastejších metód šírenia škodlivého softvéru prostredníctvom
e-mailovej komunikácie \cite{verma2017phishing}.

Významným faktorom pri takýchto útokoch prostredníctvom e-mailu
je aj využívanie sociálneho inžinierstva. Útočníci často vytvárajú
správy, ktoré pôsobia dôveryhodne alebo naliehavo, napríklad
ako faktúry, systémové upozornenia alebo bezpečnostné
aktualizácie. Cieľom takýchto techník je prinútiť používateľa,
aby vykonal požadovanú akciu bez dôkladného overenia
dôveryhodnosti správy \cite{aurangzeb2017ransomware}. 

\subsection{Mechanizmy ochrany elektronickej pošty}

\subsubsection*{Sender Policy Framework (SPF)}

je autentifikačný mechanizmus
určený na detekciu a prevenciu spoofingu e-mailových adries.
Ide o validačný protokol, ktorý umožňuje prijímajúcemu
mailovému serveru overiť, či e-mail pochádza zo servera,
ktorý je autorizovaný vlastníkmi danej domény
\cite{karim2019spam}.

Princíp SPF spočíva v publikovaní SPF záznamu v systéme
DNS danej domény. Tento záznam obsahuje zoznam IP adries
alebo serverov, ktoré sú oprávnené odosielať e-mailové
správy v mene danej domény. Pri prijatí e-mailu prijímajúci
server porovná IP adresu odosielajúceho servera so záznamom
SPF uloženým v DNS. Ak sa IP adresa v zozname nenachádza,
správa môže byť označená ako podozrivá alebo odmietnutá
\cite{rfc7208}.

SPF sa používa predovšetkým na overenie informácií v SMTP
obálke a hlavičkách správy a predstavuje jeden z najrozšírenejších
mechanizmov ochrany proti spoofingu e-mailových adries\cite{karim2019spam}.

\subsubsection*{DomainKeys Identified Mail (DKIM)}
 je kryptografický
mechanizmus autentifikácie e-mailových správ, ktorý využíva
digitálne podpisy na overenie integrity a pôvodu správy.
DKIM umožňuje prijímajúcemu serveru overiť, že obsah
správy nebol počas prenosu zmenený a že správa bola
skutočne odoslaná z deklarovanej domény
\cite{rfc6376}.

Mechanizmus DKIM využíva asymetrickú kryptografiu
založenú na dvojici kľúčov – verejnom a súkromnom kľúči.
Odosielajúci server vytvorí digitálny podpis e-mailovej
správy pomocou súkromného kľúča a tento podpis je následne
pridaný do hlavičky správy. Prijímajúci server si následne
z DNS záznamu domény odosielateľa vyžiada príslušný
verejný kľúč a pomocou neho overí platnosť podpisu
\cite{rfc6376}.

Ak je podpis platný, prijímajúci server môže potvrdiť, že
správa pochádza z deklarovanej domény a jej obsah nebol
počas prenosu modifikovaný. DKIM tak výrazne zvyšuje
dôveryhodnosť e-mailovej komunikácie a často sa používa
v kombinácii s mechanizmom SPF.

Napriek svojim výhodám DKIM nedokáže úplne zabrániť
phishingovým útokom, pretože útočníci môžu registrovať
vlastné domény a odosielať správy s platným podpisom.
Z tohto dôvodu sa DKIM zvyčajne používa spolu s ďalšími
mechanizmami ochrany elektronickej pošty
\cite{karim2019spam}.\\


\textbf{Blacklistové a reputačné systémy}

Okrem autentifikačných mechanizmov, ako sú SPF alebo DKIM,
sa pri ochrane elektronickej pošty využívajú aj blacklistové
a whitelistové systémy. Tieto mechanizmy umožňujú identifikovať
nedôveryhodných alebo naopak dôveryhodných odosielateľov
na základe zoznamov IP adries alebo domén, ktoré sú
centrálne spravované a pravidelne aktualizované.

\subsubsection*{DNS blacklist}

predstavuje databázu IP adries alebo
domén, ktoré boli identifikované ako zdroj spamových správ.
Pri tomto mechanizme e-mailový server vykoná dodatočný
DNS dotaz na špecializovaný blacklistový server s cieľom
overiť reputáciu odosielateľa. Ak sa IP adresa alebo doména
odosielateľa nachádza v blacklist databáze, e-mailová správa
môže byť automaticky odmietnutá alebo označená ako spam
\cite{karim2019spam}.

DNS blacklisty môžu fungovať dvoma hlavnými spôsobmi.
Prvý spôsob spočíva v udržiavaní databázy IP adries
mailových serverov, ktoré boli identifikované ako zdroj
spamových správ alebo ako servery zneužívané na ich
distribúciu \cite{karim2019spam}. Druhý prístup využíva
identifikáciu spamových zdrojov na základe URI adries,
najčastejšie domén alebo webových stránok, ktoré sú
používané v spamových kampaniach \cite{karim2019spam}.

Napriek svojej efektívnosti majú blacklistové systémy aj
niekoľko obmedzení. Jedným z hlavných problémov je
schopnosť útočníkov často meniť zdrojové IP adresy alebo
domény, z ktorých sú spamové správy odosielané.
Útočníci môžu napríklad využívať kompromitované
zariadenia alebo botnety, čím výrazne sťažujú identifikáciu
skutočného zdroja útoku \cite{karim2019spam}. Ďalším
problémom je čas potrebný na aktualizáciu blacklist
databáz. V prípade phishingových kampaní, ktoré využívajú
krátkodobo existujúce domény, môže byť reakčný čas
blacklist systémov nedostatočný \cite{karim2019spam}.

\subsubsection*{Whitelist}

predstavuje zoznam dôveryhodných mailových
serverov alebo odosielateľov, ktorí sú považovaní za
legitímne zdroje e-mailovej komunikácie. Správy
odoslané z týchto serverov môžu byť doručované bez
prísnej kontroly spamovými filtrami, čo znižuje
pravdepodobnosť nesprávnej klasifikácie legitímnych
správ ako spam.

Whitelisty sú často spravované priamo organizáciami,
ktoré do nich zaraďujú napríklad vlastné mailové
servery, obchodných partnerov alebo overené služby.
Používanie whitelistov môže zlepšiť efektivitu
spamových filtrov a znížiť počet falošne pozitívnych
detekcií \cite{karim2019spam}.

Významným projektom v oblasti správy blacklistov
a whitelistov je napríklad Spamhaus Project, ktorý
vznikol v roku 1998. Tento projekt poskytuje databázy
IP adries a domén spojených so spamovými aktivitami,
ktoré využívajú mnohé internetové služby a e-mailové
servery na znižovanie množstva spamových správ
doručených používateľom \cite{karim2019spam}.

\subsubsection*{Heuristické a pravidlové filtre}

 predstavujú tradičný prístup k detekcii
nevyžiadaných e-mailových správ. Tieto systémy sú založené
na súbore vopred definovaných pravidiel, ktoré analyzujú
jednotlivé časti e-mailovej správy, ako sú hlavička, predmet
správy alebo samotný obsah e-mailu. Na základe výsledku
týchto pravidiel je správe priradené určité skóre, ktoré sa
porovnáva s definovanou prahovou hodnotou. Ak celkové
skóre prekročí stanovený limit, správa je klasifikovaná ako
spam \cite{karim2019spam}.

\subsubsection*{Regex filtering.}
Jednou z najčastejších implementácií heuristických filtrov je
filtrovanie založené na regulárnych výrazoch (Regular
Expressions – regex). Regulárne výrazy umožňujú definovať
vzory textu, ktoré môžu indikovať prítomnosť spamového
obsahu. Takéto pravidlá môžu napríklad identifikovať
typické spamové frázy ako \textit{“selected winner”},
 alebo \textit{“awarded lotteries”}.
Ak e-mail obsahuje takéto výrazy, systém mu priradí určitú
váhu a pri prekročení definovaného prahu môže byť správa
označená ako spam \cite{jakobsson2007phishing}.

\subsubsection*{Heuristic filtering.}
Heuristické filtre analyzujú aj hlavičku e-mailovej správy,
ktorá obsahuje informácie o pôvode, obsahu a prenose
správy. Spammeri často manipulujú tieto údaje pomocou
techník ako spoofing, aby správa pôsobila legitímne.
Heuristické filtre dokážu takéto nezrovnalosti identifikovať
napríklad kontrolou formátu hlavičky alebo analýzou
reťazca serverov v záznamoch \textit{Received}
\cite{jakobsson2007phishing,karim2019spam}.

\subsubsection*{Content-based filtre}

predstavujú metódu detekcie spamových
správ založenú na analýze samotného obsahu e-mailovej
správy. Tento prístup skúma text správy, predmet e-mailu
alebo prípadné prílohy s cieľom identifikovať znaky typické
pre spamové kampane \cite{karim2019spam}.

Pri tomto prístupe sa jednotlivým slovám alebo frázam
priraďuje určitá váha podľa toho, ako často sa vyskytujú
v spamových správach. Následne sa vypočíta celkové skóre
správy, ktoré sa porovnáva s vopred definovanou prahovou
hodnotou. Ak výsledné skóre prekročí túto hodnotu,
e-mailová správa je klasifikovaná ako spam
\cite{karim2019spam}.

Typickým príkladom sú pravidlové expertové systémy,
ktoré využívajú zoznam charakteristických spamových
kľúčových slov. Napríklad slová ako \textit{cheap},
\textit{mortgage}  môžu byť
považované za indikátory spamovej komunikácie a
prispievajú k výslednému skóre správy
\cite{karim2019spam}.

Nevýhodou content-based filtrov je citlivosť na kontext.
Systém môže označiť legitímny e-mail ako spam,
ak obsahuje slová typické pre spamové správy,
hoci sú použité v inom kontexte. Tento problém je
známy ako \textit{false positive} detekcia.

Content-based filtre sa od jednoduchých heuristických filtrov odlišujú tým,
že sa zameriavajú na celkový význam a štatistické vlastnosti obsahu správy,
nielen na prítomnosť izolovaných zakázaných výrazov.
\subsubsection*{Kolaboratívne systémy}

Kolaboratívne spamové filtre využívajú informácie získané
od viacerých používateľov alebo e-mailových serverov.
Základnou myšlienkou je, že spamové správy sú často
rozosielané vo veľkom množstve identických alebo veľmi
podobných e-mailov. Ak je určitá správa identifikovaná
ako spam jedným používateľom alebo serverom, táto
informácia môže byť zdieľaná s ostatnými systémami
\cite{karim2019spam}.

\subsection*{Hash sharing.}
Pri tejto metóde sa z obsahu e-mailovej správy vypočíta
kryptografický hash, ktorý slúži ako jedinečný identifikátor
správy. Hash je alfanumerický reťazec, ktorý reprezentuje
digitálny odtlačok správy. Poskytovatelia e-mailových služieb
môžu tieto hodnoty ukladať do databáz a porovnávať ich
s hash hodnotami nových prichádzajúcich správ. Ak sa hash
hodnota zhoduje s hodnotou uloženou v databáze spamových
správ, e-mail môže byť okamžite označený ako spam
\cite{karim2019spam}.

\subsubsection*{Distributed Checksum Clearinghouse (DCC)}predstavuje
kolaboratívny systém na detekciu spamových správ.
Systém vypočíta kontrolný súčet e-mailovej správy
a ukladá ho do distribuovanej databázy. Pri každom
ďalšom označení správy ako spam sa počet výskytov
daného kontrolného súčtu zvyšuje. Ak je určitý kontrolný
súčet zaznamenaný veľkým počtom používateľov,
systém môže s vysokou pravdepodobnosťou identifikovať
takúto správu ako spamovú kampaň \cite{jakobsson2007phishing,karim2019spam}.

Nevýhodou týchto systémov je, že aj malá zmena v obsahu správy môže
zmeniť výslednú hash hodnotu, čo útočníci využívajú pri obchádzaní
takýchto mechanizmov.

\textbf{Moderné metódy detekcie}

S rastúcim objemom elektronickej komunikácie a čoraz sofistikovanejšími technikami spamových kampaní sa tradičné metódy filtrovania e-mailov, založené výlučne na statických pravidlách alebo blacklistoch, ukázali ako nedostatočné. Moderné systémy ochrany elektronickej pošty preto čoraz častejšie využívajú prístupy založené na strojovom učení (machine learning) a hlbokom učení (deep learning), ktoré umožňujú automaticky identifikovať vzory typické pre spamové správy a adaptovať sa na nové typy útokov \cite{karim2019spam}).

\subsubsection*{Machine learning prístupy}
predstavujú významný pokrok v oblasti detekcie spamu, pretože umožňujú automatickú klasifikáciu e-mailových správ na základe analýzy rôznych charakteristík správy. Takéto systémy sú založené na trénovaní klasifikačných modelov na označených datasetoch, ktoré obsahujú príklady legitímnych e-mailov (ham) a spamových správ. Model následne dokáže identifikovať pravdepodobnosť, že nová správa patrí do jednej z týchto kategórií\cite{karim2019spam}.

Pri klasifikácii e-mailov sa využívajú rôzne typy príznakov, medzi ktoré patria napríklad obsah správy, frekvencia výskytu slov, štruktúra e-mailu, prítomnosť URL adries alebo informácie z hlavičky správy. Medzi najčastejšie používané algoritmy strojového učenia patria napríklad Support Vector Machines, Random Forest alebo logistická regresia. Tieto modely dokážu identifikovať komplexné vzory v dátach a dosahujú vysokú presnosť pri detekcii spamových správ \cite{karim2019spam}.


\subsubsection*{Naive Bayes klasifikátor}

je jednou z najpoužívanejších metód strojového učenia pri detekcii spamu je Naive Bayes klasifikátor. Tento algoritmus je založený na Bayesovej pravdepodobnostnej teórii a predpoklade nezávislosti jednotlivých príznakov. V kontexte filtrovania e-mailov sa Naive Bayes používa na výpočet pravdepodobnosti, že daná správa je spam na základe výskytu určitých slov alebo charakteristík v texte e-mailu\cite{karim2019spam}.

Model je trénovaný na veľkom množstve e-mailových správ, z ktorých sa vypočítava pravdepodobnosť výskytu jednotlivých slov v spamových a legitímnych správach. Pri analýze novej správy algoritmus vypočíta pravdepodobnosť jej zaradenia do každej kategórie a následne ju klasifikuje ako spam alebo legitímny e-mail. Výhodou tejto metódy je jej relatívna jednoduchosť, nízka výpočtová náročnosť a vysoká efektivita pri spracovaní veľkého množstva e-mailov\cite{karim2019spam}.

\subsubsection*{Deep learning prístupy.}

V posledných rokoch sa v oblasti detekcie spamových e-mailov čoraz viac uplatňujú aj metódy hlbokého učenia. Deep learning využíva neurónové siete s viacerými vrstvami, ktoré dokážu automaticky extrahovať komplexné reprezentácie dát bez potreby manuálneho definovania príznakov. Tieto modely dokážu analyzovať nielen jednotlivé slová, ale aj kontext a vzťahy medzi nimi\cite{karim2019spam}.

Medzi často používané architektúry patria napríklad konvolučné neurónové siete (CNN) alebo rekurentné neurónové siete (RNN), ktoré sú vhodné na spracovanie textových dát. Deep learning modely dokážu identifikovať komplexné jazykové vzory a štruktúry typické pre spamové správy, čím zvyšujú presnosť detekcie v porovnaní s tradičnými metódami. Napriek tomu však vyžadujú väčšie množstvo tréningových dát a vyšší výpočtový výkon \cite{karim2019spam}.

V súčasnosti sa moderné systémy ochrany elektronickej pošty často opierajú o hybridné riešenia, ktoré kombinujú tradičné mechanizmy, ako sú blacklisty alebo heuristické pravidlá, s pokročilými metódami strojového učenia a hlbokého učenia. Takýto viacvrstvový prístup umožňuje efektívnejšie identifikovať spamové správy a reagovať na nové formy e-mailových hrozieb\cite{karim2019spam}.


\subsection{Analýza existujúcich nástrojov}

\textbf{MailHog} je open-source nástroj určený predovšetkým na testovanie emailových aplikácií počas vývoja. Zachytáva odosielané emaily pomocou lokálneho SMTP servera a umožňuje ich zobrazenie prostredníctvom webového rozhrania. Vývojári tak môžu testovať odosielanie emailov bez rizika, že budú skutočne doručené externým používateľom \cite{mailhog}.

\textbf{GoPhish} je open-source framework určený na simuláciu phishingových kampaní. Umožňuje vytvárať phishingové emaily, pristávacie stránky a sledovať interakcie používateľov, napríklad kliknutia na odkazy alebo zadanie prihlasovacích údajov. Nástroj sa často využíva pri bezpečnostných školeniach a testovaní povedomia používateľov o phishingových útokoch \cite{gophish}.

\textbf{PhishTool} je platforma zameraná na analýzu phishingových emailov a podporu reakcie na bezpečnostné incidenty. Umožňuje bezpečnostným analytikom zhromažďovať, analyzovať a kategorizovať phishingové správy. Systém poskytuje nástroje na extrakciu indikátorov kompromitácie (IOC), analýzu hlavičiek emailov a spoluprácu medzi členmi bezpečnostného tímu \cite{phishtool}.

Napriek tomu PhishTool neposkytuje automatizovanú klasifikáciu alebo rozhodovanie na základe kombinácie heuristických a modelových prístupov. Systém funguje primárne ako analytická pomôcka, ktorá vizualizuje jednotlivé indikátory a ponecháva finálne rozhodnutie na používateľovi. Identifikované nezrovnalosti sú prezentované izolovane bez komplexného skórovania alebo agregovaného vyhodnotenia rizika.

Na Obr.~\ref{fig:phishtool_auth} je zobrazený príklad autentifikačnej analýzy, kde všetky mechanizmy (SPF, DKIM, DMARC) dosahujú stav \textit{PASS}. Takýto výsledok môže viesť k dojmu legitimity správy, aj keď iné indikátory môžu naznačovať potenciálne riziko.

\begin{figure}[H]
\centering
\includegraphics[width=\columnwidth]{phishtool_auth.png}
\caption{Autentifikačné výsledky (SPF, DKIM, DMARC) v nástroji PhishTool}
\label{fig:phishtool_auth}
\end{figure}

Na Obr.~\ref{fig:phishtool_display} je zobrazená detekcia nekonzistentného zobrazovaného mena (\textit{display-name}), kde nástroj upozorňuje na nesúlad medzi názvom odosielateľa a emailovou adresou. Tento indikátor je prezentovaný ako samostatné upozornenie bez ďalšieho kontextu alebo kvantifikácie rizika.

\begin{figure}[H]
\centering
\includegraphics[width=\columnwidth]{phishtool_display.png}
\caption{Detekcia nekonzistentného display-name}
\label{fig:phishtool_display}
\end{figure}

Obr.~\ref{fig:phishtool_returnpath} ukazuje nesúlad medzi doménou \texttt{Return-Path} a doménou odosielateľa (\texttt{From}). Aj keď ide o relevantný indikátor potenciálneho spoofingu, PhishTool ho opäť prezentuje len ako izolované upozornenie.

\begin{figure}[H]
\centering
\includegraphics[width=\columnwidth]{phishtool_returnpath.png}
\caption{Nesúlad domény Return-Path a From}
\label{fig:phishtool_returnpath}
\end{figure}

Tieto príklady ilustrujú hlavné obmedzenie nástroja: jednotlivé indikátory sú síce detegované korektne, avšak nie sú agregované do jednotného rozhodnutia ani skóre rizika. Na rozdiel od navrhnutého riešenia tak chýba schopnosť automaticky kombinovať viacero signálov a vyhodnotiť ich v kontexte, čo je kľúčové pre robustnú detekciu phishingových kampaní.

\paragraph{Príklad nalýza konkrétneho e-mailu}

Na základe analyzovaného e-mailu (pozri \ref{fig:phishtool_auth}--\ref{fig:phishtool_returnpath}) je možné detailne ilustrovať správanie nástroja PhishTool a zároveň poukázať na limity čisto indikátorového prístupu.

Z pohľadu autentifikačných mechanizmov e-mail vykazuje konzistentné výsledky:
SPF, DKIM aj DMARC validácie sú úspešné (\textit{PASS}), čo znamená,
že správa bola odoslaná zo serverov autorizovaných doménou
\texttt{packeta.sk} a jej obsah nebol po ceste modifikovaný.
Takýto výsledok je typický pre legitímne transakčné e-maily a
predstavuje silný signál dôveryhodnosti.

Napriek tomu PhishTool identifikuje viaceré čiastkové nezrovnalosti.
Konkrétne ide o:
a) nekonzistentný display-name,
b) nesúlad medzi doménou \texttt{Return-Path} (\texttt{mg.packeta.sk})
a doménou odosielateľa (\texttt{packeta.sk}).

Tieto indikátory sú však v tomto prípade vysvetliteľné legitímnym spôsobom.
Použitie subdomény \texttt{mg.packeta.sk} súvisí s využitím služby
Mailgun na hromadné odosielanie e-mailov, kde je bežné používať
odlišný \texttt{Return-Path} pre správu bounce správ.
Rovnako display-name „Packeta Slovakia“ je konzistentný s doménou
odosielateľa a nepredstavuje reálny pokus o spoofing.

Obsahová analýza e-mailu taktiež podporuje jeho legitimitu.
Správa obsahuje konkrétne informácie o zásielke (identifikátor,
miesto doručenia, stav „zaplatené“) a neobsahuje typické znaky
phishingu ako výzvy na zadanie prihlasovacích údajov, urgentné
výzvy na akciu alebo podozrivé odkazy.
Text je konzistentný, jazykovo korektný a zodpovedá štandardným
notifikačným e-mailom logistických služieb. :contentReference[oaicite:0]{index=0}

Z pohľadu celkového vyhodnotenia ide teda o \textbf{legitímny transakčný e-mail},
ktorý však obsahuje niektoré technické anomálie bežné v moderných
mailingových infraštruktúrach.

Tento príklad poukazuje na zásadný problém izolovaného heuristického
prístupu: jednotlivé indikátory môžu signalizovať potenciálne riziko,
avšak bez kontextu vedú k falošným pozitívam. PhishTool tieto indikátory
správne deteguje, no neposkytuje mechanizmus na ich agregáciu ani
interpretáciu v širšom kontexte.

Naopak, navrhnutý systém kombinuje heuristiky s modelom strojového učenia,
čo umožňuje zohľadniť kontext správy a správne klasifikovať takýto e-mail
ako legitímny aj napriek prítomnosti čiastkových anomálií.

\section{Experimentálna časť}
\subsection{Návrh experimentálneho prostredia a zberu emailových dát}
Pre potreby experimentálnej časti práce bolo navrhnuté a implementované lokálne testovacie prostredie na platforme Windows s využitím nástroja Docker Compose. Cieľom bolo vytvoriť reprodukovateľné prostredie, v ktorom bude možné simulovať phishingové kampane, zachytávať doručené správy a následne ich analyzovať heuristickými aj modelovými metódami.

V prostredí boli nasadené dve hlavné služby. Prvou službou bol MailHog, ktorý slúžil ako lokálny SMTP server a zároveň poskytoval webové používateľské rozhranie a API na čítanie prijatých správ. MailHog bol sprístupnený na portoch 1025 pre SMTP komunikáciu a 8025 pre webové rozhranie. Druhou službou bol GoPhish, ktorý bol použitý na správu a spúšťanie testovacích phishingových kampaní. GoPhish bol nakonfigurovaný s administrátorským rozhraním na porte 3333 a phishing serverom na porte 8080.

V rámci laboratórneho nasadenia bola v GoPhish použitá lokálna konfigurácia bez TLS, konkrétne s nastaveniami ADMIN\_USE\_TLS=false a PHISH\_USE\_TLS=false. Tento prístup bol zvolený zámerne, keďže cieľom nebolo nasadenie do produkčného prostredia, ale jednoduché a transparentné testovanie v izolovanom lokálnom prostredí. SMTP profil v GoPhish bol nastavený na adresu mailhog:1025 pod profilom MailHog Local, čím bolo zabezpečené, že všetky odoslané správy z kampaní budú doručené do MailHog a budú dostupné na ďalšie spracovanie.

\textbf{Príprava vstupných datasetov}

Ďalším krokom bolo zostavenie vhodného dátového základu pre tréning a testovanie navrhovaných klasifikačných modelov. Vstupné datasety boli získané z verejne dostupných zdrojov na platforme Kaggle. Použité boli viaceré emailové kolekcie, ktoré sa líšili pôvodom, štruktúrou aj typom anotácie. Medzi hlavné zdroje patrili napríklad Enron, SpamAssassin, CEAS\_08, phishing\_email, Nazario, Nigerian\_Fraud a ďalšie.

Už po stiahnutí a základnej kontrole dát sa ukázalo, že jednotlivé datasety sú výrazne nevyvážené a zároveň heterogénne z pohľadu obsahu aj označenia tried. Niektoré datasety obsahovali kombináciu legitímnych a škodlivých správ označených binárnym alebo viackategóriovým labelom, zatiaľ čo iné boli špecializované len na konkrétny typ správ. Napríklad dataset \textit{Nazario} obsahoval prakticky výlučne phishingové správy a dataset \textit{Nigerian\_Fraud} bol zameraný na správy reprezentujúce finančné podvody.

Významný problém predstavovalo aj to, že zastúpenie jednotlivých zdrojov v celkovom datasete nebolo rovnomerné. Dominantnú časť tvorili správy z datasetu Enron, zatiaľ čo viaceré ostatné zdroje boli zastúpené len v malom percente. Takáto nerovnováha mohla viesť k tomu, že model by sa počas tréningu učil skôr špecifické štýlové alebo zdrojové charakteristiky konkrétnych datasetov než všeobecné znaky phishingových a podvodných emailov.

Percentuálne zastúpenie jednotlivých zdrojov po stiahnutí datasetov je znázornené na obrázku~\ref{fig:download_dataset_distribution}. Z grafu je zrejmé, že dataset Enron tvorí dominantnú časť dátovej základne, pričom ostatné datasety majú výrazne menší podiel.
\begin{figure}[H]
    \centering
    \includegraphics[width=1\linewidth]{download_dataset_distribution_pie.pdf}
    \caption{Percentuálne zastúpenie jednotlivých zdrojových datasetov po stiahnutí}
    \label{fig:download_dataset_distribution}
\end{figure}


\textbf{Normalizácia a unifikácia dátovej schémy}

Na základe tejto analýzy z predchadzajúcej kapitoly bolo potrebné v ďalších realizovať čistenie dát, zjednotenie schémy, deduplikáciu a návrh takých tréningových scenárov, ktoré by obmedzili vplyv nerovnováhy a zdrojového skreslenia na výsledný model.

Na tento účel bol najprv implementovaný skript. Úlohou skriptu bolo načítať vstupné CSV súbory z rôznych datasetov, vyčistiť textové polia, odstrániť neplatné a duplicitné záznamy a zjednotiť výsledok do spoločnej schémy.

Každý email bol reprezentovaný textovými poľami subject, body a text, pričom sa zachovávala aj surová verzia textu (subject\_raw, body\_raw, text\_raw) pre potreby neskoršieho auditu. Tento dvojitý zápis – surový text a modelový text – sa ukázal ako dôležitý z hľadiska transparentnosti, pretože umožňoval spätne skontrolovať, akým spôsobom boli texty normalizované a ktoré transformácie boli aplikované pred tréningom modelu.

Normalizácia zahŕňala najmä:
\begin{itemize}
    \item odstránenie nadbytočných whitespace znakov,
    \item zjednotenie koncov riadkov,
    \item odstránenie technického šumu z hlavičiek vložených do tela správ,
    \item nahradenie URL, emailových adries a čísel špeciálnymi tokenmi, napríklad <URL>, <EMAIL> a <NUM>,
    \item základné čistenie textu od neštandardných znakov.
\end{itemize}

Týmto spôsobom sa znižovalo riziko, že model sa naučí nerelevantné povrchové vzory viazané na konkrétne adresy, domény alebo identifikátory správ. Súčasne sa vytvoril priestor pre robustnejšie porovnanie modelov naprieč rôznymi scenármi.

V počiatočnej fáze práce bol použitý klasický prístup: zlúčenie viacerých verejných emailových zdrojov, čistenie, deduplikácia a následné rozdelenie na tréningový, validačný a testovací split. Tento prístup bol metodicky funkčný a viedol k vysokým metrikám na štandardnom teste, avšak ďalšie analýzy ukázali, že takéto výsledky môžu byť čiastočne nafúknuté prítomnosťou skratkovitých signálov, ktoré nesúvisia priamo so semantikou emailu.

Ako kritický problém sa ukázalo najmä source leakage a template leakage. V praxi to znamená, že model sa mohol naučiť rozpoznávať zdroj datasetu alebo typickú šablónu kampane namiesto všeobecných znakov phishingu, spamu alebo finančného podvodu. Indikáciou tohto problému boli:
\begin{itemize}
    \item veľmi vysoké výsledky na jednoduchom IID teste,
    \item výrazné rozdiely pri testovaní na robustnejších splitoch,
    \item vysoká source purity niektorých zdrojov,
    \item dobré výsledky aj pri baseline modeloch, ktoré mali prístup len k metadátam alebo informácii o zdroji.
\end{itemize}

Na elimináciu týchto problémov bola pipeline postupne rozšírená o viacero obranných mechanizmov. Najprv bola zavedená silnejšia textová normalizácia, následne source-aware spracovanie, explicitné limitovanie dominantných zdrojov, tokenové maskovanie výrazov prezrádzajúcich pôvod datasetu a napokon aj template-aware splitovanie a template-aware regulácia tréningových scenárov.

Výsledkom nebol jeden jediný dataset, ale sada tréningových scenárov A–M, ktoré zdieľali spoločné evaluačné splity a líšili sa iba stratégiou konštrukcie tréningovej množiny. Takýto návrh umožnil systematicky skúmať, ktoré opatrenia majú pozitívny vplyv na robustnosť modelu a ktoré naopak vedú len k optickému zlepšeniu metrík.

\textbf{Návrh spoločných evaluačných splitov}

Jednou z najdôležitejších metodických zmien bolo zavedenie zdieľaných splitov, ktoré sú spoločné pre všetky scenáre. Vďaka tomu bolo možné férovo porovnávať jednotlivé tréningové stratégie, keďže rozdiely vo výsledkoch nevznikali z iného náhodného rozdelenia dát, ale priamo zo spôsobu konštrukcie tréningovej množiny.

Takto navrhnuté rozdelenie má viacero výhod. Sedemdesiatpäť percent dát v train\_pool poskytuje dostatok variability na učenie rozhodovacích hraníc aj pri minoritných triedach. Desaťpercentný validačný split je dostatočne veľký na stabilný model selection, no zároveň neuberá príliš veľa tréningových dát. Desaťpercentný IID test zabezpečuje porovnateľnosť s bežnou literatúrou, zatiaľ čo deployment a hard splity poskytujú realistickejší obraz o generalizácii mimo komfortnej distribúcie.
\begin{table}[H]
\centering
\caption{Rozdelenie datasetu na tréningové, validačné a testovacie splity}
\label{tab:split_design}
\footnotesize
\setlength{\tabcolsep}{3pt}
\renewcommand{\arraystretch}{1.1}
\begin{tabular}{|l|p{4.8cm}|c|}
\hline
\textbf{Split} & \textbf{Charakteristika} & \textbf{Pomer} \\
\hline
train\_pool & Základná množina pre tréning modelov a tvorbu scenárov & 0.75 \\
val\_iid & Validačný split na výber modelu a ladenie hyperparametrov & 0.10 \\
test\_iid & Test so zachovanou distribúciou dát (IID) pre benchmark & 0.10 \\
test\_deployment & Test simulujúci reálne nasadenie modelu & 0.05 \\
test\_hard\_source & Test robustnosti pri zmene zdrojového datasetu & 0.02 \\
test\_hard\_cluster & Test robustnosti pri zmene šablón/clusterov & 0.02 \\
\hline
\end{tabular}
\end{table}

Osobitný význam majú splity test\_hard\_source a test\_hard\_cluster. Tie fungujú ako stress test modelu. Ich cieľom nie je nahradiť hlavný benchmark, ale odhaliť, či model ostáva funkčný aj pri zmene zdroja alebo pri nových šablónach emailov. Práve tieto splity umožnili lepšie odlíšiť modely, ktoré sa učia všeobecné znaky útoku, od modelov, ktoré sa spoliehajú na datasetové artefakty.

\textbf{Evolúcia scenárov A--M}
\begin{table*}[!b]
\centering
\caption{Evolúcia scenárov A--M: mechanizmus, prínos a hlavné obmedzenia}
\label{tab:scenarios_a_m}
\footnotesize
\setlength{\tabcolsep}{4pt}
\renewcommand{\arraystretch}{1.1}
\begin{tabular}{|c|p{4.2cm}|p{5.2cm}|p{5.2cm}|}
\hline
\textbf{Scenár} & \textbf{Čo sa menilo} & \textbf{Čo fungovalo} & \textbf{Hlavný problém / riziko} \\
\hline
A & Strict balanced train (rovnaký počet vzoriek na triedu) & Stabilný tréning, vyvážené triedy & Umelá distribúcia; nízka podobnosť s reálnym nasadením \\
B & Full train + implicitné váhovanie tried & Lepšie využitie dát, silný baseline & Dominancia zdrojov; pretrváva source leakage \\
C & Capped train (obmedzenie veľkých tried) & Zníženie prevahy majority tried & Nerieši source ani template leakage \\
D & Oversampled balanced train & Lepší recall minoritných tried & Overfitting na duplicity a syntetické vzory \\
E & Syntetická augmentácia vybraných tried & Pomoc pre minority triedy & Riziko synthetic signature \\
F & Mild rebalance + LLM augmentácia & Väčšia textová variabilita & Nekonzistentná kvalita generovaných vzoriek \\
G & Source-balanced train (cap podľa zdrojov) & Zníženie dominancie veľkých zdrojov & Časť source shortcutov zostáva \\
H & Anti-source token sanitization & Redukcia explicitných dataset markerov & Nerieši template leakage \\
I & Strict source cap & Lepšia robustnosť voči source shiftu & Strata časti kvalitných dát \\
J & Limit dominancie zdroja v triede & Silný baseline na hard splitoch & Model sa môže učiť template vzory \\
K & Drop high-leak sources & Výrazné potlačenie leakage & Nižšia variabilita dát \\
L & Kombinovaný anti-source (K+J+I+H) + hard-legit injection & Silný anti-shortcut baseline & Pretrváva template leakage a feature shortcuty \\
M & L + anti-template regulácia + feature dropout + synthetic mix & Najvyššia robustnosť (source + template) & Mierny pokles surových metrík \\
\hline
\end{tabular}
\end{table*}
Jednotlivé experimentálne scenáre boli navrhnuté ako postupná evolúcia od základných sampling stratégií až po metodicky prísne anti-leak konfigurácie. Ich cieľom nebolo iba maximalizovať klasické metriky (napr. accuracy), ale predovšetkým analyzovať, ako rôzne zásahy do dát ovplyvňujú robustnosť modelu a jeho schopnosť generalizovať mimo tréningovej distribúcie.

Pre zachovanie férového porovnania všetky scenáre využívali rovnaké validačné a testovacie splity; líšili sa výlučne zložením tréningovej množiny. Prehľad jednotlivých scenárov vrátane ich mechanizmu, prínosu a hlavných obmedzení je uvedený v tabuľke~\ref{tab:scenarios_a_m}.



Scenáre A--D predstavujú základné prístupy k riešeniu nerovnováhy tried v dátach. Zahŕňajú striktne vyvážený tréningový set, využitie celého datasetu s implicitným váhovaním, obmedzenie veľkých tried (capped train) a oversampling minoritných tried. Tieto prístupy umožnili analyzovať vplyv distribúcie tried na výkon modelu, avšak samy osebe neriešili problém tzv. dataset leakage, najmä závislosti od zdroja alebo šablóny emailu.

Scenáre E a F rozširujú pipeline o syntetickú augmentáciu dát, vrátane využitia LLM modelov. Tento prístup priniesol zvýšenie variability textov a zlepšenie výsledkov na minoritných triedach. Na druhej strane sa ukázalo, že model sa môže naučiť rozpoznávať charakteristické znaky generovaných dát (tzv. synthetic signature), čo znižuje jeho schopnosť generalizácie na reálne dáta. Z tohto dôvodu boli tieto scenáre vnímané skôr ako experimentálne než ako finálne riešenie.

Nasledujúce scenáre (G--K) sa zameriavajú na elimináciu source leakage. Scenár G zavádza vyváženie podľa zdrojov, čím sa znižuje dominancia veľkých datasetov. Scenár H dopĺňa textovú normalizáciu (token sanitization), ktorá odstraňuje explicitné znaky prezrádzajúce pôvod dát. Scenáre I a J ďalej sprísňujú obmedzenia na zastúpenie jednotlivých zdrojov, pričom scenár J zavádza limit dominancie zdroja v rámci triedy. Scenár K ide ešte ďalej a úplne odstraňuje vybrané high-leak zdroje, ktoré vykazovali silné datasetové artefakty.

Scenár L predstavuje kombinovaný anti-source prístup, ktorý integruje viacero predchádzajúcich techník (K+J+I+H).Tento scenár sa ukázal ako veľmi silný baseline z hľadiska robustnosti voči source-based shortcutom.

Ďalšia analýza však odhalila, že aj pri eliminácii source leakage môže model využívať tzv. template leakage, teda učiť sa špecifické štruktúry alebo šablóny emailových kampaní. Zároveň sa ukázala potenciálna závislosť na niektorých ručne navrhnutých (engineered) príznakoch, najmä z emailových hlavičiek.

Z tohto dôvodu bol navrhnutý finálny scenár M, označený ako \textit{anti-template feature-regularized}. Tento scenár rozširuje anti-leak prístup o ďalšie mechanizmy:

\begin{itemize}
    \item limit počtu vzoriek na template (cluster\_id) v rámci každej triedy,
    \item obmedzenie kombinácie template + source,
    \item label-aware syntetickú augmentáciu s kontrolovaným podielom,
    \item feature dropout na príznakoch náchylných na shortcuty.
\end{itemize}

Variant M predstavuje metodicky najprísnejší prístup, pretože cielene obmedzuje nielen source leakage, ale aj template leakage a závislosť modelu od povrchových technických znakov. Aj keď môže dochádzať k miernemu poklesu niektorých surových metrík, tento scenár poskytuje najlepší predpoklad pre robustnú generalizáciu modelu v reálnych podmienkach.

\textbf{Feature engineering a heuristické signály}

Okrem textovej reprezentácie emailov boli extrahované aj základné číselné príznaky založené na hlavičkách a obsahu správ. Zahŕňali napríklad charakteristiky URL (počet odkazov, skrátené URL, unikátne domény), nesúlad medzi identitami odosielateľa (\textit{From}, \textit{Reply-To}, \textit{Return-Path}), výskyt podozrivých TLD a indikátory typosquattingu. Súčasťou boli aj jednoduché obsahové signály, ako počty kľúčových slov (urgentné, autentifikačné, finančné) a štatistiky textu (pomer veľkých písmen, číslic, výkričníkov).

Tieto príznaky boli využité jednak v heuristickom skórovaní, jednak ako doplnkové vstupy v hybridných modeloch. Z dôvodu obmedzenia závislosti modelu od týchto príznakov bol v pokročilých scenároch aplikovaný prístup feature regularization a čiastočný feature dropout.

\textbf{Tréning modelov a porovnávané architektúry}

Na klasifikáciu emailov boli použité predovšetkým lineárne modely vhodné pre sparse textové reprezentácie. Ich hlavnými výhodami sú rýchlosť tréningu, nízka pamäťová náročnosť, stabilita a dobrá interpretovateľnosť. V experimentoch boli porovnávané najmä nasledujúce architektúry:

\begin{itemize}
    \item hybrid\_logreg,
    \item hybrid\_linear\_svc\_cal,
    \item hybrid\_sgd\_log,
    \item hybrid\_sgd\_hinge.
\end{itemize}

\textbf{Výber vstupnej reprezentácie}

Kľúčovým rozhodnutím pri návrhu klasifikačného modelu bolo určenie vhodnej vstupnej reprezentácie emailových správ. V počiatočných experimentoch boli analyzované redukované pohľady na dáta, konkrétne \texttt{header\_only}, \texttt{subject\_only} a \texttt{body\_only}. Tieto varianty umožnili pochopiť, aký typ informácie nesie najväčší diskriminačný signál.

Finálne experimenty však ukázali, že najlepší výkon poskytuje kombinovaný vstup \texttt{text\_plus\_features}, ktorý spája textový obsah správy (predmet a telo emailu) s odvodenými číselnými príznakmi založenými na hlavičkách a štruktúre správy. Tento prístup umožňuje modelu využívať nielen sémantický obsah, ale aj doplnkové heuristické signály relevantné z pohľadu emailovej bezpečnosti.

Porovnanie jednotlivých vstupných reprezentácií pre finálny model scenára M je znázornené na obrázku~\ref{fig:input_comparison_final_m}. Z grafu je zrejmé, že varianty \texttt{header\_only} a \texttt{subject\_only} dosahujú výrazne nižšie hodnoty macro F1, čo naznačuje, že tieto časti emailu samy osebe neposkytujú dostatočne robustný signál pre klasifikáciu.

Reprezentácia \texttt{body\_only} prináša výrazné zlepšenie, čím sa potvrdzuje, že rozhodujúca informácia je obsiahnutá najmä v texte správy. Napriek tomu však najlepšie výsledky dosahuje kombinovaný prístup \texttt{text\_plus\_features}, ktorý konzistentne prekonáva všetky ostatné varianty naprieč evaluačnými splitmi vrátane náročných scenárov \texttt{test\_hard\_source}, \texttt{test\_hard\_cluster} a \texttt{test\_deployment}.

Tento výsledok potvrdzuje, že vhodne navrhnuté heuristické príznaky dokážu efektívne doplniť textový model bez toho, aby dominovali jeho rozhodovaniu. Kombinácia obsahových a štrukturálnych signálov tak predstavuje najrobustnejší základ pre klasifikáciu emailov v podmienkach blízkych reálnemu nasadeniu.

\begin{figure}[H]
    \centering
    \includegraphics[width=1\linewidth]{final_model_m_input_representation.pdf}
    \caption{Porovnanie vstupných reprezentácií pre finálny model scenára M (hybrid\_logreg).}
    \label{fig:input_comparison_final_m}
\end{figure}

\textbf{Single-stage vs. two-stage prístup} 

V experimentálnej časti boli porovnané dva základné prístupy ku klasifikácii. Prvým bol single-stage model, ktorý priamo priraďuje email do jednej z cieľových tried. Druhým bol two-stage prístup, pri ktorom prvá fáza rozhoduje medzi benign a suspicious správami a druhá fáza následne rozlišuje konkrétnu kategóriu medzi podozrivými správami.

Two-stage prístup je intuitívne atraktívny, pretože kopíruje viacúrovňové rozhodovanie známe z bezpečnostných systémov. V našich experimentoch však nepriniesol konzistentné zlepšenie oproti single-stage klasifikácii. Naopak, vo viacerých benchmarkoch boli single-stage modely stabilnejšie a jednoduchšie interpretovateľné. Okrem toho majú single-stage modely nižšiu implementačnú a prevádzkovú komplexitu, čo je dôležité z pohľadu budúceho nasadenia.

Z týchto dôvodov bol single-stage prístup zvolený ako hlavný produkčný smer a two-stage bol ponechaný ako porovnávacia vetva vhodná do experimentálnej a diskusnej časti práce.

\textbf{Evaluačný protokol a výber metrík}

Evaluačný protokol bol navrhnutý tak, aby striktne oddeľoval fázu výberu modelu od finálneho hodnotenia a zároveň umožňoval systematické testovanie robustnosti pri distribučných posunoch. Výber modelu (model selection) prebiehal výhradne na validačnom splite \textit{val\_iid}. Hlavný benchmark bol následne reportovaný na \textit{test\_iid} a finálne modely boli ďalej hodnotené na robustnostných splitoch \textit{test\_hard\_source}, \textit{test\_hard\_cluster} a \textit{test\_deployment}.

Ako hlavná metrika bola zvolená \textbf{macro F1}, ktorá v multiclass nastavení poskytuje vyvážené hodnotenie naprieč všetkými triedami, vrátane minoritných kategórií. Okrem nej boli sledované aj doplnkové metriky:

\begin{itemize}
    \item balanced\_accuracy,
    \item accuracy,
    \item per-class precision, recall a F1,
    \item confusion matrix.
\end{itemize}

V počiatočnej binárnej fáze bola hlavná pozornosť venovaná metrikám \textbf{precision}, \textbf{recall} a \textbf{F1} pre pozitívnu triedu (phishing). Metrika accuracy bola interpretovaná len ako doplnková, keďže pri nevyvážených dátach môže viesť k skreslenému hodnoteniu výkonu modelu. Tento prístup umožnil lepšie zachytiť schopnosť modelu detegovať phishingové správy pri zachovaní kontroly nad počtom falošných pozitívnych a negatívnych predikcií.

\subsubsection{Výber finálnych scenárov a modelov}

Na základe realizovaných experimentov sa ako najrelevantnejší ukázal scenár M a jeho tréningové varianty, ktoré počas experimentov ďalej evolvovali. Pôvodný scenár M poskytol metodicky čistý anti-template základ, avšak počas tréningu sa ukázalo, že cielená úprava distribúcie minoritných tried vedie k ďalšiemu zlepšeniu multiclass výkonu na kampanových dátach.

Z tohto dôvodu bol ako hlavný model zvolený tréningový variant:
\begin{itemize}
    \item \textbf{Primárny model:} \\
    \textit{scenario\_m\_rebalanced+ hybrid\_logreg + text\_plus\_features}
\end{itemize}

Tento variant vznikol z pôvodného scenára M cieleným rebalancingom minoritných tried, najmä \texttt{financial\_fraud} a \texttt{spam}, a následne bol ďalej jemne doladený triedovým škálovaním.

Ako sekundárny model bol ponechaný metodicky čistý variant pôvodného scenára M:
\begin{itemize}
    \item \textbf{Sekundárny model:} \\
    \textit{scenario\_m\_anti\_template\_feature\_regularized + hybrid\_logreg + text\_plus\_features}
\end{itemize}

Pre referenčné porovnanie ostal zachovaný aj silný anti-source baseline zo scenára L:
\begin{itemize}
    \item \textbf{Referenčný baseline:} \\
    \textit{scenario\_l\_combined\_anti\_source + hybrid\_sgd\_log + text\_plus\_features}
\end{itemize}

Takto definovaný výber lepšie odráža reálny tréningový vývoj než pôvodné porovnanie L verzus M v izolovanej podobe.

\subsubsection{Porovnanie finálnych kandidátov na robustnostných splitoch}

Táto časť vychádza výhradne z hodnotenia na spoločných validačných a testovacích splitoch datasetu
(\textit{val\_iid}, \textit{test\_iid}, \textit{test\_hard\_source}, \textit{test\_hard\_cluster}, \textit{test\_deployment}).


Na náročných splitoch boli dosiahnuté tieto hodnoty \texttt{f1\_macro}:

\begin{table}[H]
\centering
\caption{Porovnanie výsledkov modelov na rôznych testovacích scenároch}
\label{tab:model_results}
\footnotesize
\setlength{\tabcolsep}{5pt}
\renewcommand{\arraystretch}{1.1}
\begin{tabular}{lccc}
\toprule
\textbf{Dataset} & \textbf{L+SGD} & \textbf{M+SGD} & \textbf{M+LogReg} \\
\midrule
test\_hard\_source   & 0.694 & 0.801 & 0.802 \\
test\_hard\_cluster  & 0.820 & 0.811 & 0.848 \\
test\_deployment     & 0.876 & 0.895 & 0.896 \\
\bottomrule
\end{tabular}

\end{table}

Výsledky potvrdzujú, že scenár M je robustnejší na source-shift splite
a konfigurácia \textit{M + hybrid\_logreg} dosahuje najlepší celkový kompromis na náročných testoch.

\begin{figure}[H]
    \centering
    \includegraphics[width=1\linewidth]{lm_hard_grouped_models.pdf}
    \caption{Porovnanie finálnych kandidátov na hard a deployment splitoch (macro F1)}
    \label{fig:lm_hard_grouped}
\end{figure}


\textbf{Analýza finálneho modelu}

Ako finálny model bol zvolený variant:
\texttt{scenario\_m\_rebalanced\_fraud\_spam + hybrid\_logreg + text\_plus\_features}
s jemným triedovým škálovaním pre \texttt{financial\_fraud:1.1}.

Pre objektívne zhodnotenie bol model analyzovaný na nezávislom testovacom splite \texttt{test\_deployment}, ktorý lepšie reprezentuje podmienky reálneho nasadenia než experimentálne dáta z kampane.

Konfúzna matica finálneho modelu je znázornená na obrázku~\ref{fig:confusion_matrix}. Z výsledkov je možné pozorovať niekoľko kľúčových charakteristík.

Trieda \texttt{legit} je klasifikovaná veľmi spoľahlivo, čo dokumentuje vysoký počet správne klasifikovaných vzoriek (14748) a relatívne nízky počet chýb smerom k ostatným triedam (najmä 244 do \texttt{phishing} a 245 do \texttt{spam}). To naznačuje, že model dokáže robustne odlišovať bežnú komunikáciu od škodlivého obsahu.

Trieda \texttt{phishing} dosahuje taktiež veľmi dobré výsledky (2642 správne klasifikovaných vzoriek), pričom chyby sú rozdelené medzi \texttt{spam} (46 prípadov) a \texttt{legit} (36 prípadov). Tento vzor chýb odráža prirodzené prekrytie medzi manipulatívnym marketingovým obsahom a phishingovými technikami.

Trieda \texttt{financial\_fraud} je klasifikovaná stabilne (336 správnych klasifikácií), pričom väčšina chýb smeruje do triedy \texttt{spam} (20 prípadov). Tento jav je očakávaný, keďže finančné podvody často využívajú podobné jazykové vzory ako agresívny marketing alebo spam.

Najproblematickejšou triedou ostáva \texttt{spam}. Aj keď model správne klasifikoval 805 vzoriek, dochádza k zámene so \texttt{phishing} (11 prípadov) a \texttt{legit} (13 prípadov). To naznačuje, že hranica medzi spamom a ostatnými kategóriami je v reálnych dátach menej jednoznačná a výrazne závisí od kontextu správy.

\begin{figure}[H]
    \centering
    \includegraphics[width=1\linewidth]{final_model_confusion_matrix_test_deployment.pdf}
    \caption{Konfúzna matica finálneho modelu na splite test\_deployment}
    \label{fig:confusion_matrix}
\end{figure}

Z pohľadu praktického nasadenia model vykazuje veľmi dobrú schopnosť rozlišovať medzi legitímnou a škodlivou komunikáciou. Najväčšie chyby vznikajú pri jemnej diferenciácii medzi jednotlivými typmi škodlivých správ (spam vs. phishing vs. financial fraud), čo je typické aj pre reálne bezpečnostné systémy.

Dôležitým pozorovaním je rozdiel medzi experimentálnymi dátami (Experiment 1) a testovacím splitom. Kým v Experiment 1 modely dosahovali takmer nulový počet false positive prípadov, na \texttt{test\_deployment} sa objavuje realistickejší trade-off medzi precision a recall. Tento rozdiel potvrdzuje, že model nie je preučený na syntetické alebo kampanové dáta, ale zachováva schopnosť generalizácie na širšiu distribúciu emailov.

Z pohľadu metodickej čistoty aj praktického použitia sa tento variant javí ako najvhodnejší finálny model. Sekundárny model založený na pôvodnom scenári M ostáva referenčným bodom a baseline zo scenára L slúži ako porovnávacia anti-source konfigurácia.

\subsubsection{Zhrnutie navrhnutého riešenia}

Navrhnuté riešenie prešlo vývojom od pôvodnej jednoduchej pipeline založenej na zlúčení a náhodnom rozdelení verejných datasetov až po robustný, reprodukovateľný a metodicky kontrolovaný experimentálny rámec. Bolo vytvorené lokálne laboratórne prostredie s nástrojmi GoPhish a MailHog, pripravený automatizovaný pipeline pre sťahovanie, čistenie a unifikáciu datasetov, implementované mechanizmy proti source a template leakage a navrhnutý viacúrovňový evaluačný protokol.

Hlavným prínosom navrhnutého riešenia nie je len samotná klasifikácia emailov, ale predovšetkým dôraz na to, aby sa model učil robustné a všeobecne prenositeľné signály. Z tohto dôvodu bola značná pozornosť venovaná split stratégii, anti-leak scenárom, auditovateľnosti dátovej prípravy a porovnaniu viacerých alternatívnych tréningových konfigurácií.

Takto pripravený rámec vytvára pevný základ pre nasledujúcu experimentálnu časť práce, v ktorej budú systematicky porovnané heuristické, modelové a hybridné prístupy na klasifikáciu emailových správ.


\subsection{Heuristický analyzátor}
\begin{table*}[!t]
\label{tab:heuristic_scoring_rules}
\centering
\caption{Pravidlá heuristického skórovania analyzátora}
\footnotesize
\renewcommand{\arraystretch}{1.15}

\begin{tabular}{p{4.5cm} p{2.8cm} p{9cm}}
\toprule
\textbf{Indikátor} & \textbf{Príspevok do skóre} & \textbf{Význam} \\
\midrule
SPF fail & +20 & Silný technický signál zlyhania autentifikácie odosielateľa \\
DKIM fail & +15 & Indikuje problém s integritou alebo podpisom správy \\
DMARC fail & +20 & Potvrdzuje nesúlad domény a autentifikačných mechanizmov \\
Nesúlad From a Reply-To & +15 & Častý znak spoofingu alebo presmerovania odpovede \\
Nesúlad From a Return-Path & +15 & Indikuje rozdiel medzi deklarovaným a technickým odosielateľom \\
IP adresa v URL & +10 za URL, max. +20 & URL smerujúca priamo na IP býva častá pri phishingu \\
Skracovač URL & +5 za URL, max. +10 & Skrýva cieľovú doménu a znižuje transparentnosť odkazu \\
Obfuskovaná URL & +3 za URL, max. +10 & URL s @, percent-encodingom alebo dlhým query reťazcom môže skrývať skutočný cieľ \\
Anchor mismatch & +5 za výskyt, max. +10 & Rozdiel medzi zobrazeným odkazom a skutočným cieľom \\
Podozrivá TLD & +5 za doménu, max. +10 & Výskyt TLD častejšie spájaných s phishingom alebo spamom \\
Externé URL domény & +3 za doménu, max. +12 & Odkazy vedúce mimo domény odosielateľa \\
Vysoký nesúlad odosielateľ / URL domény & +5 až +10 & Silný indikátor, že obsah nesedí s identitou odosielateľa \\
Brand typosquatting & +15 & Doména napodobňuje známu značku s malou zmenou názvu \\
Credential keywords & +3 za výskyt, max. +15 & Výrazy ako login, password, account, confirm \\
Urgent keywords & +2 za výskyt, max. +10 & Výrazy vytvárajúce tlak a naliehavosť \\
Vysoká hustota credential / urgent výrazov & +4 až +5 & Koncentrovaný výskyt manipulatívnych výrazov v texte \\
Riziková príloha & +10 & Príloha s potenciálne nebezpečnou príponou \\
Display-name spoofing & +10 & Zobrazené meno dôveryhodnej značky pri cudzej doméne \\
Unicode mixed script & +10 & Kombinácia skriptov v doméne môže indikovať homografický útok \\
Received anomália & +5 & Neštandardný transportný reťazec správy \\
Vysoký pomer veľkých písmen & +5 & Agresívny alebo manipulatívny štýl komunikácie \\
Veľa výkričníkov & +5 & Typický znak nátlakového alebo spamového štýlu \\
Trusted domain guardrail & $-20$ & Zníženie skóre pri dôveryhodnej komunikácii bez silných rizikových signálov \\
\bottomrule
\end{tabular}
\end{table*}
\textbf{Architektúra a motivácia}

Heuristický analyzátor bol navrhnutý ako transparentná a auditovateľná vrstva spracovania e-mailových správ, ktorej úlohou je extrahovať technické aj obsahové príznaky zo správ vo formáte \texttt{.eml}, kvantifikovať ich pomocou explicitne definovaného skórovacieho mechanizmu a pripraviť ich ako vstup pre ďalšie modelové spracovanie. Motiváciou pre jeho návrh bolo, že samotný text správy neposkytuje dostatočné informácie na spoľahlivú detekciu útokov, keďže významná časť indikátorov je obsiahnutá v hlavičkách, URL adresách a štruktúre správy. Zároveň bolo cieľom zachovať vysokú mieru interpretovateľnosti, aby bolo možné spätne vysvetliť rozhodnutia systému.

\textbf{Parsovanie a reprezentácia správy}

Spracovanie začína parsovaním \texttt{.eml} súboru pomocou triedy \texttt{BytesParser} z modulu \texttt{email}, čím sa zachováva plná MIME štruktúra správy. Analyzátor extrahuje hlavičky, textové časti, HTML časti a prílohy, pričom každá kategória je spracovaná samostatne. V prípade multipart správ sa prechádza celý MIME strom a jednotlivé časti sa klasifikujú podľa typu obsahu a atribútov. Výsledkom je jednotná reprezentácia obsahujúca textové segmenty, HTML obsah, prílohy a zoznam URL adries, ktorá slúži ako základ pre ďalšiu extrakciu príznakov.

\textbf{Extrakcia URL a doménových príznakov}

URL adresy sú detegované kombináciou základného regulárneho výrazu a pokročilého vzoru podporujúceho Unicode a IDN domény. Každá URL je normalizovaná, aby sa odstránili redundancie a zabezpečila konzistentná reprezentácia. Následne sa analyzuje hostiteľ URL, pričom sa sleduje výskyt IP adries, skracovačov URL, podozrivých TLD domén, punycode zápisu a zmiešaných Unicode skriptov. Dôležitým príznakom je aj pomer URL domén, ktoré sa nezhodujú s doménou odosielateľa, čo indikuje potenciálne presmerovanie na nedôveryhodné zdroje. Samostatne sa deteguje aj typosquatting značiek pomocou Levenshteinovej vzdialenosti voči známym doménam.

\textbf{Hlavičkové a autentifikačné indikátory}

Analyzátor extrahuje kľúčové hlavičky ako \texttt{From}, \texttt{Reply-To}, \texttt{Return-Path} a \texttt{Message-ID} a porovnáva ich domény s cieľom identifikovať nesúlady. Tieto nesúlady predstavujú častý znak spoofingu. Okrem toho sa spracúva hlavička \texttt{Authentication-Results}, z ktorej sa extrahujú výsledky SPF, DKIM a DMARC. Zlyhania týchto mechanizmov sú považované za silné technické indikátory podozrivej správy, keďže signalizujú problém s autenticitou odosielateľa.

\textbf{Obsahové a štylistické príznaky}

Text správy je analyzovaný z hľadiska štylistických znakov a výskytu kľúčových slov. Sleduje sa počet urgentných, credential, threat a finančných výrazov, ako aj ich hustota v texte. Tieto príznaky reflektujú typické stratégie sociálneho inžinierstva, ktoré využívajú tlak, naliehavosť alebo manipuláciu používateľa. Okrem toho sa počíta pomer veľkých písmen, počet výkričníkov a pomer číslic, ktoré môžu indikovať agresívny alebo manipulatívny štýl komunikácie.

\textbf{HTML a štrukturálne anomálie}

V HTML častiach sa detegujú nesúlady medzi zobrazeným textom odkazu a jeho skutočným cieľom, čo predstavuje typickú phishingovú techniku. Analyzátor tiež identifikuje obfuskované URL adresy, ktoré obsahujú netypické znaky alebo nadmernú komplexitu. Sleduje sa aj prítomnosť HTML-only správ a anomálie v počte \texttt{Received} hlavičiek, ktoré môžu indikovať neštandardný transport správy.

\textbf{Heuristické skórovanie}

Na základe extrahovaných príznakov analyzátor vypočíta heuristické skóre v rozsahu 0 až 100. Každý indikátor prispieva k skóre definovaným počtom bodov, pričom silné technické signály ako zlyhanie SPF, DKIM alebo DMARC majú vyššiu váhu než slabšie obsahové indikátory. Skóre je následne mapované na tri úrovne rizika: nízke, stredné a vysoké. Súčasťou výstupu je aj zoznam dôvodov (\texttt{heuristic\_reasons}), ktorý zaznamenáva, ktoré indikátory prispeli k výslednému skóre, čím sa zabezpečuje vysvetliteľnosť systému.

\textbf{Ochrana proti falošným poplachom}

Analyzátor obsahuje aj mechanizmus na zníženie počtu false positive prípadov. Ak správa pochádza z dôveryhodnej domény, autentifikačné mechanizmy neindikujú problém a neobsahuje silné rizikové znaky, výsledné skóre je redukované. Tento prístup reflektuje potrebu praktickej použiteľnosti systému, kde je dôležité minimalizovať zbytočné upozornenia pri legitímnej komunikácii.


\textbf{Integrácia do pipeline}

Výstup analyzátora má dvojitú úlohu. Na jednej strane poskytuje samostatné heuristické hodnotenie s vysvetlením, ktoré môže byť využité priamo používateľom. Na druhej strane generuje štruktúrovaný dataset obsahujúci numerické a kategóriálne príznaky, ktorý slúži ako vstup pre modely strojového učenia. Tento dizajn umožňuje kombinovať výhody heuristického prístupu (interpretovateľnosť) s výhodami ML prístupov (vyššia presnosť a generalizácia).

\subsection{Ukážka výstupu analyzátora}
\begin{figure}[H]
\centering
\includegraphics[width=0.8\columnwidth]{analyzer_output.png}
\caption{Príklad výstupu analyzátora e-mailov}
\label{fig:analyzer_output}
\end{figure}


Na obr.~\ref{fig:analyzer_output} je zobrazený príklad výstupu analyzátora. Výstup je štruktúrovaný do viacerých blokov, ktoré kombinujú heuristickú analýzu a model strojového učenia. Vidíme základné metadáta e-mailu (odosielateľ, príjemca, predmet, počet URL), heuristické skóre spolu s vysvetlením (napr. prítomnosť \texttt{credential\_keywords}), zoznam detegovaných príznakov (v tomto prípade žiadne) a výstup ML modelu vrátane pravdepodobností tried.

Heuristická aj ML vrstva v danom príklade klasifikujú správu ako legitímnu, pričom ML model poskytuje detailnejší pohľad cez distribúciu pravdepodobností. Tento kombinovaný výstup umožňuje lepšiu interpretáciu rozhodnutí a identifikáciu potenciálnych rizík.

\subsection{Experimenty}

\subsubsection{Experiment 1 – Simulácia kampane a klasifikácia doručených e-mailov}\label{epx1}

Cieľom Experimentu 1 bolo overiť funkčnosť navrhnutého end-to-end postupu v kontrolovanom prostredí, teda od generovania e-mailových správ, cez ich distribúciu v simulovanej kampani, až po ich zber, analýzu a klasifikáciu. Experiment zároveň slúžil na prvé overenie správania heuristického analyzátora aj modelových klasifikátorov na správach doručených v prostredí GoPhish + MailHog.

\paragraph{Generovanie dát}

Vstupné správy boli zostavené kombináciou viacerých prístupov s cieľom dosiahnuť obsahovú aj štylistickú variabilitu. Použité boli štyri hlavné vetvy:

\begin{itemize}
    \item \textbf{manuálne vytvorené správy} – menší počet kontrolovaných emailov pripravených autorom, reprezentujúcich základné legitímne aj škodlivé scenáre,
    \item \textbf{LLM drafty} (\texttt{exp1\_llm\_draft}) – návrhy správ generované jazykovým modelom,
    \item \textbf{mutované správy} (\texttt{exp1\_mutated}) – automaticky upravené existujúce správy so zmenami v texte, URL a hlavičkách,
    \item \textbf{realistické úpravy z viacerých zdrojov} (\texttt{exp1\_multi\_real\_edit}) – manuálne alebo poloautomaticky upravené správy inšpirované reálnymi vzormi.
\end{itemize}
\paragraph{Simulácia kampane a zber dát}

Správy boli následne seedované do nástroja GoPhish a doručované cez MailHog. Na rozdiel od jednoduchého jednopríjemcového testovania boli použité \textbf{unikátne recipient identity} vo formáte \texttt{user\{id\}@local.lab}, čím sa dosiahol realistickejší zber a lepšia identifikácia jednotlivých správ v kampani. Po seedingu boli doručené správy zozbierané vo formáte \texttt{.eml} a následne analyzované skriptom \texttt{analyze\_campaign.py}, ktorý zabezpečil parsovanie správ, extrakciu feature-ov a prípravu vstupov pre klasifikáciu.

\subsubsection*{Zloženie dát}

\begin{table}[http]
\centering
\caption{Zloženie dát a rozsah spracovania v experimente~1}
\label{tab:exp1_dataset}
\footnotesize
\setlength{\tabcolsep}{5pt}
\renewcommand{\arraystretch}{1.1}
\begin{tabular}{lc}
\toprule
\textbf{Metrika} & \textbf{Hodnota} \\
\midrule
Počet seedovaných správ celkom & 250 \\
\quad vytvorené manuálne & 25 \\
\quad \texttt{exp1\_llm\_draft} & 25 \\
\quad \texttt{exp1\_mutated} & 200 \\
\quad \texttt{exp1\_multi\_real\_edit} & 25 \\
\addlinespace
Počet zozbieraných správ (aktuálny run) & 250 \\
Distribúcia binárnych labelov (0/1) & 50 / 200 \\
\bottomrule
\end{tabular}
\end{table}

\paragraph{Kategórie a mapovanie tried}

V experimente boli použité štyri cieľové kategórie:
\texttt{legit}, \texttt{phishing}, \texttt{spam} a \texttt{financial\_fraud}.  
V prípadoch, kde sa v návrhu správ objavoval scenár \texttt{spoofing}, bol na účely jednotného evaluačného prostredia mapovaný do triedy \texttt{phishing}, keďže z pohľadu klasifikačného cieľa išlo o formu podvodnej sociálnej manipulácie.

\paragraph{Heuristický analyzátor}

Heuristický analyzátor bol pôvodne nastavený s prahovou hodnotou 40. Pri tomto nastavení však vykazoval nízky recall, teda nedokázal zachytiť dostatočný počet škodlivých správ. Z tohto dôvodu bol vykonaný sweep prahových hodnôt a ako najvhodnejší kompromis pre Experiment 1 bol zvolený prah 20.

\begin{figure}[H]
    \centering
    \includegraphics[width=0.9\linewidth]{exp1_heuristic_threshold_sweep.pdf}
    \caption{Závislosť metrík heuristického analyzátora od prahovej hodnoty.}
    \label{fig:exp1_threshold_sweep}
\end{figure}

Z grafu (Obr.~\ref{fig:exp1_threshold_sweep}) je možné pozorovať typický trade-off medzi precision a recall. So znižovaním prahu dochádza k výraznému nárastu recall, avšak zároveň k poklesu precision. Zvolený prah 20 predstavuje vyvážený kompromis, ktorý umožňuje zachytiť väčšinu škodlivých správ pri stále akceptovateľnej miere falošne pozitívnych detekcií.

\paragraph{Porovnanie heuristiky a modelov}

Na rovnakých dátach bol vykonaný aj priamy porovnávací experiment medzi heuristickým analyzátorom a modelovými prístupmi. Výsledky sú znázornené na obrázku~\ref{fig:exp1_model_comparison}.

\begin{figure}[H]
    \centering
    \includegraphics[width=0.9\linewidth]{exp1_model_comparison.pdf}
    \caption{Porovnanie F1, balanced accuracy a MCC medzi heuristikou a modelmi.}
    \label{fig:exp1_model_comparison}
\end{figure}

Z grafu je zrejmé, že heuristický analyzátor dosahuje vysoký recall a F1 skóre, avšak výrazne zaostáva v metrikách balanced accuracy a MCC. Naopak, modelové prístupy poskytujú stabilnejšiu separáciu tried a lepší kompromis medzi precision a recall :contentReference[oaicite:1]{index=1}.

\paragraph{Multiclass modelové výsledky scenárov L a M}

Na dátach z aktuálneho sender-clean runu boli následne vyhodnotené scenáre L a M. Vývoj multiclass výkonu je znázornený grafom na obrázku~\ref{fig:exp1_multiclass_comparison}.

\begin{figure}[H]
    \centering
    \includegraphics[width=0.9\linewidth]{exp1_multiclass_method_comparison.pdf}
    \caption{Porovnanie multiclass výkonu heuristiky, modelu a ich kombinácie.}
    \label{fig:exp1_multiclass_comparison}
\end{figure}

Výsledky ukazujú, že samotná heuristika dosahuje výrazne nižší multiclass výkon než modelové prístupy. Modely založené na scenári M dosahujú najlepšie výsledky, pričom kombinácia textu a feature-ov (\texttt{text\_plus\_features}) prináša výrazné zlepšenie oproti čisto textovému vstupu :contentReference[oaicite:2]{index=2}.

\paragraph{Analýza finálneho modelu na dátach experimentu}

Detailnejší pohľad na správanie finálneho modelu poskytuje konfúzna matica na obrázku~\ref{fig:exp1_confusion_matrix}.

\begin{figure}[H]
    \centering
    \includegraphics[width=0.9\linewidth]{exp1_m_confusion_best_final.pdf}
    \caption{Konfúzna matica finálneho modelu na dátach Experimentu 1.}
    \label{fig:exp1_confusion_matrix}
\end{figure}

Z konfúznej matice je možné pozorovať, že model veľmi spoľahlivo klasifikuje triedy \texttt{legit} a \texttt{financial\_fraud}, pričom najväčšie problémy vznikajú pri triede \texttt{spam}. Tá je často zamieňaná s \texttt{phishing}.

\paragraph{Interpretácia výsledkov}

Experiment potvrdil niekoľko kľúčových zistení:

\begin{itemize}
    \item kvalita seed vstupu (najmä konzistencia subject/body a profil odosielateľa) má zásadný vplyv na výsledky,
    \item heuristický prístup je vhodný ako citlivý screening (vyšší recall),
    \item modelové prístupy poskytujú lepšiu separáciu tried a stabilnejšie metriky,
    \item kombinácia \texttt{text\_plus\_features} je výrazne silnejšia než samotný text,
    \item rebalancing tried v scenári M výrazne pomáha minoritným kategóriám.
\end{itemize}

\paragraph{Záver experimentu}

Experiment 1 potvrdil funkčnosť celej pipeline:
\[
\text{generovanie} \rightarrow \text{seeding} \rightarrow \text{zber} \rightarrow \text{analýza} \rightarrow \text{klasifikácia}.
\]

Z pohľadu výsledkov sa ako najlepší model pre tento typ experimentov ukázal variant:

\texttt{scenario\_m\_rebalanced\_fraud\_spam + hybrid\_logreg + text\_plus\_features}


Tento model poskytol najlepší kompromis medzi multiclass výkonom, robustnosťou a schopnosťou generalizácie na reálne kampanové dáta. V porovnaní s heuristikou dosahuje stabilnejšie výsledky a lepšiu separáciu tried, pričom zároveň eliminuje závislosť na jednoduchých pravidlových signáloch.

Heuristický analyzátor ostáva vhodným doplnkom pre prvotný screening, avšak pre finálnu klasifikáciu a rozhodovanie je model založený na scenári M jednoznačne vhodnejšou voľbou.

\subsubsection{Experiment 2 – Analýza hlavičiek a autentifikačných indikátorov}\label{epx2}

Cieľom Experimentu 2 bolo overiť dve komplementárne vlastnosti navrhnutého riešenia. 
Prvou bolo zistiť, či heuristický analyzátor spoľahlivo deteguje vybrané hlavičkové a autentifikačné indikátory v kontrolovaných podmienkach. Druhou bolo overiť, ako sa tieto indikátory, samostatne aj v kombinácii s modelom strojového učenia, správajú na realistickejších kampanových dátach generovaných v prostredí GoPhish a MailHog.

Dôležitým aspektom tohto experimentu je, že \textbf{heuristický analyzátor je navrhnutý ako binárny detektor}, ktorý rozhoduje medzi triedami \texttt{legit} a \texttt{suspicious}. Rovnaký princíp bol použitý aj pri hodnotení modelu strojového učenia, kde multiclass výstup bol interpretovaný aj v binárnom pohľade (legit vs. akýkoľvek typ škodlivej správy). Multiclass metriky preto slúžia skôr ako doplnkové hodnotenie správania modelu.

\paragraph{Analyzované heuristické indikátory}

V rámci Experimentu 2 boli definované a testované konkrétne heuristické indikátory zamerané na analýzu e-mailových hlavičiek a autentifikačných mechanizmov. Tieto indikátory vychádzajú priamo z cieľov práce, najmä detekcie spoofingu, nesúladu hlavičiek a overovania autentifikačných kontrol.

Medzi sledované indikátory patrili:

\begin{itemize}
    \item \texttt{from\_reply\_mismatch} – nesúlad medzi adresami \texttt{From} a \texttt{Reply-To},
    \item \texttt{from\_return\_path\_mismatch} – nesúlad medzi \texttt{From} a \texttt{Return-Path},
    \item \texttt{display\_name\_spoof\_flag} – podozrivý zobrazovaný názov odosielateľa (napr. imitácia známej služby),
    \item \texttt{message\_id\_domain\_mismatch} – nesúlad domény v hlavičke \texttt{Message-ID},
    \item \texttt{received\_anomaly} – neštandardná alebo podozrivá reťaz hlavičiek \texttt{Received},
    \item \texttt{spf\_fail} – zlyhanie SPF autentifikácie,
    \item \texttt{dkim\_fail} – zlyhanie DKIM podpisu,
    \item \texttt{dmarc\_fail} – zlyhanie DMARC politiky,
    \item \texttt{trusted\_domain\_only} – príznak legitímnej komunikácie pochádzajúcej výlučne z dôveryhodnej domény.
\end{itemize}

Tieto indikátory boli cielene manipulované v kontrolovanom datasete (Vetva A), pričom pre každý prípad bol definovaný očakávaný výstup.


\paragraph{Vetva A – kontrolovaný header/auth dataset}

Prvá vetva bola zameraná na presné overenie detekcie jednotlivých indikátorov. 
Bol vytvorený kontrolovaný súbor \texttt{.eml} správ, v ktorom boli cielene manipulované vybrané položky hlavičky.

Pre každý prípad bol definovaný očakávaný výstup a následne porovnaný s detekciou analyzátora. Výsledky ukázali, že pre všetky sledované indikátory dosiahol analyzátor perfektné výsledky (accuracy, precision aj recall rovné 1.0).

Tento výsledok potvrdzuje, že heuristický analyzátor spoľahlivo identifikuje explicitne definované hlavičkové a autentifikačné anomálie. V tejto vetve však nejde o klasifikačný problém, ale o validáciu detekčných pravidiel. Z tohoto dôvodu je taká vysoká uspešnosť.

\paragraph{Vetva B – realistická kampaň}

Druhá vetva experimentu bola realizovaná ako simulovaná phishingová kampaň v prostredí GoPhish + MailHog. Dataset obsahoval 100 správ so zmiešaným pôvodom (LLM generované a realisticky upravené správy) a štyrmi triedami: \texttt{legit}, \texttt{phishing}, \texttt{spam} a \texttt{financial\_fraud}.

Na týchto dátach boli porovnané tri prístupy:

\begin{itemize}
    \item \textbf{heuristic\_only} – binárne rozhodovanie legit vs suspicious,
    \item \textbf{ml\_only} – multiclass model interpretovaný aj v binárnom režime,
    \item \textbf{heuristic\_plus\_ml} – kombinovaný prístup.
\end{itemize}

\paragraph{Porovnanie prístupov}

Pri hodnotení výsledkov je potrebné zdôrazniť, že primárnym cieľom heuristického analyzátora je binárne rozhodovanie medzi triedami \texttt{legit} a \texttt{suspicious}. Z tohto dôvodu sú hlavné metriky reportované v binárnom režime, zatiaľ čo multiclass výsledky slúžia len ako doplnkové hodnotenie správania modelov.

V binárnom hodnotení (legit vs. suspicious) dosiahli jednotlivé prístupy nasledovné výsledky:

\begin{table}[H]
\centering
\caption{Porovnanie prístupov v Experimente 2 (Vetva B)}
\label{tab:exp2_comparison}
\footnotesize
\setlength{\tabcolsep}{4pt}
\renewcommand{\arraystretch}{1.1}
\begin{tabular}{lccccc}
\toprule
\textbf{Prístup} & \textbf{Acc} & \textbf{Bal Acc} & \textbf{Prec} & \textbf{Rec} & \textbf{F1} \\
\midrule
\texttt{heuristic\_only}        & 0.80 & 0.50  & 0.80   & 1.00 & 0.8889 \\
\texttt{ml\_only}               & \textbf{0.95} & \textbf{0.875} & \textbf{0.9412} & 1.00 & \textbf{0.9697} \\
\texttt{heuristic\_plus\_ml}    & 0.85 & 0.625 & 0.8421 & 1.00 & 0.9143 \\
\bottomrule
\end{tabular}
\end{table}

Výsledky potvrdzujú, že heuristický analyzátor je veľmi citlivý na podozrivé správy, čo sa prejavuje maximálnym recall (\texttt{1.00}). Na druhej strane má nižšiu špecificitu, čo vedie k horšej hodnote balanced accuracy (\texttt{0.50}). Inými slovami, heuristika má tendenciu označovať väčší počet správ ako podozrivé, aj keď časť z nich patrí do triedy \texttt{legit}.

Najlepší výkon dosiahol prístup \texttt{ml\_only}, ktorý poskytuje výrazne lepší kompromis medzi precision a recall a zároveň dosahuje najvyššiu hodnotu F1 skóre (\texttt{0.9697}). Kombinovaný prístup \texttt{heuristic\_plus\_ml} v tomto nastavení nepriniesol zlepšenie, čo naznačuje, že jednoduché prepojenie heuristiky a modelu bez dodatočného ladenia nemusí byť optimálne.

Z doplnkového pohľadu boli analyzované aj multiclass výsledky, ktoré však nie sú primárnym cieľom tohto experimentu. V tomto režime heuristika dosahuje nízke hodnoty (napr. \texttt{mc\_f1\_macro = 0.13}), čo je očakávané, keďže nie je navrhnutá na jemné rozlišovanie medzi triedami \texttt{phishing}, \texttt{spam} a \texttt{financial\_fraud}. Model strojového učenia dosahuje v multiclass nastavení výrazne lepšie výsledky (\texttt{mc\_f1\_macro = 0.53}), čo potvrdzuje jeho schopnosť detailnejšej kategorizácie.

Tieto výsledky ukazujú, že heuristika je vhodná najmä ako detektor podozrivosti (binárna úloha), zatiaľ čo model strojového učenia je nevyhnutný pre presnú multiclass klasifikáciu.

\paragraph{Interpretácia výsledkov}

Výsledky experimentu poukazujú na rozdielnu úlohu jednotlivých prístupov:

\begin{itemize}
    \item heuristický analyzátor funguje spoľahlivo ako detektor prítomnosti rizikových indikátorov,
    \item model strojového učenia poskytuje lepšiu schopnosť rozlišovať medzi jednotlivými typmi útokov,
    \item priame kombinovanie oboch prístupov bez dodatočného ladenia môže viesť k degradácii výkonu.
\end{itemize}

Z pohľadu návrhu systému sa preto ako vhodné riešenie javí architektúra, v ktorej heuristický analyzátor slúži ako doplnkový alebo signalizačný modul (napr. pre vysvetliteľnosť alebo zvýraznenie rizikových príznakov), zatiaľ čo finálne rozhodovanie o kategórii správy je ponechané na model strojového učenia. Táto architektúra však nebola v rámci Experimentu 2 explicitne implementovaná ako produkčný pipeline, ale vyplýva z interpretácie dosiahnutých výsledkov.

\paragraph{Záver experimentu}

Experiment 2 potvrdil, že heuristický analyzátor a model strojového učenia majú v systéme odlišné, ale komplementárne úlohy.

Heuristický analyzátor spoľahlivo deteguje explicitné bezpečnostné indikátory a funguje ako prvá vrstva rozhodovania (legit vs suspicious). Model strojového učenia následne zabezpečuje jemnejšie rozlíšenie medzi jednotlivými typmi škodlivých správ.

Z metodického hľadiska sa preto ako najvhodnejší prístup javí využitie heuristiky ako doplnkového a vysvetľujúceho modulu, zatiaľ čo hlavné rozhodovanie o kategórii správy by malo byť ponechané na model strojového učenia.


\subsubsection{Experiment 3 – URL indikátory a ich vplyv na detekciu podozrivých správ}\label{epx3}

Cieľom Experimentu 3 bolo overiť správanie URL-orientovaných heuristických indikátorov a zistiť, do akej miery tieto signály prispievajú k rozlišovaniu medzi legitímnymi a podozrivými e-mailami. Experiment bol rozdelený na dve vetvy: kontrolovaný benchmark (Track A) a realistickú URL-heavy kampaň (Track B).

\paragraph{Track A – validácia URL indikátorov}

V kontrolovanom prostredí boli všetky sledované URL indikátory detegované bez chyby (accuracy, precision aj recall rovné 1.0). Tento výsledok potvrdzuje správnosť implementácie detekčných pravidiel a spoľahlivú identifikáciu jednotlivých URL znakov.

Pri binárnom rozhodovaní heuristika dosiahla nižšie hodnoty (accuracy = 0.78), čo ukazuje, že kombinovanie indikátorov do finálneho rozhodnutia je náročnejšie ako samotná detekcia jednotlivých znakov.

Je však potrebné zdôrazniť, že ide o kontrolovaný dataset, a preto tieto výsledky predstavujú validáciu implementácie detekčných pravidiel, nie hodnotenie správania analyzátora v realistickom prostredí.

\paragraph{Track B – realistická URL-heavy kampaň}

Na realistickej kampani boli porovnané tri prístupy: \texttt{heuristic\_only}, \texttt{ml\_only} a \texttt{heuristic\_plus\_ml}.

\begin{table}[h]
\centering
\caption{Porovnanie prístupov v binárnom hodnotení (Experiment 3 – Track B)}
\begin{tabular}{lcccc}
\hline
Metóda & Accuracy & Precision & Recall & F1 \\
\hline
heuristic\_only & 0.80 & 0.80 & 1.00 & 0.8889 \\
ml\_only & 0.94 & 0.9625 & 0.9625 & 0.9625 \\
heuristic\_plus\_ml & 0.87 & 0.8602 & 1.00 & 0.9249 \\
\hline
\end{tabular}
\end{table}

Heuristický analyzátor dosahuje maximálny recall, čo znamená, že zachytáva všetky podozrivé správy. Na druhej strane však nedokáže efektívne odlišovať legitímne správy, čo sa prejavuje nulovým počtom true negative prípadov a vyšším počtom false positive.

Model strojového učenia dosahuje najlepší celkový výkon a poskytuje vyvážený kompromis medzi precision a recall. Kombinovaný prístup zlepšuje heuristiku, avšak nedosahuje výkon samostatného modelu, čo naznačuje, že jednoduché pravidlá kombinácie nemusia byť optimálne.

\paragraph{Analýza URL indikátorov}

V kampani sa jednotlivé URL indikátory vyskytovali rovnomerne (približne 25 výskytov na indikátor), čo zabezpečuje vyvážené testovacie prostredie. Analýza korelácie ukázala, že nie všetky indikátory majú rovnaký vplyv na heuristické skóre.


\begin{figure}[H]
    \centering
    \includegraphics[width=0.9\linewidth]{exp3_url_indicator_spearman_correlation.pdf}
    \caption{Spearmanova korelácia medzi prítomnosťou URL indikátorov a heuristickým skóre v Experimente 3 (Track B). Najsilnejší pozitívny vzťah bol pozorovaný pri indikátore \texttt{brand\_typosquat}, zatiaľ čo \texttt{shortener\_url} a \texttt{suspicious\_tld} vykázali slabší až negatívny vzťah. Indikátor \texttt{anchor\_mismatch} sa v tejto vetve experimentu neobjavil.}
    \label{fig:exp3_url_indicator_correlation}
\end{figure}

Doplnková analýza korelácie ukázala, že jednotlivé URL indikátory neprispievajú k heuristickému skóre rovnako. Najsilnejší pozitívny vzťah bol pozorovaný pri indikátore \texttt{brand\_typosquat}, zatiaľ čo indikátory \texttt{shortener\_url} a \texttt{suspicious\_tld} vykazovali v tomto datasete slabší, resp. negatívny vzťah ku skóre. Tento výsledok naznačuje, že diskriminačná sila jednotlivých URL príznakov závisí aj od konkrétneho zloženia kampanových dát.

\paragraph{Analýza chýb}

Heuristický analyzátor vykazuje zvýšenú chybovosť najmä pri absencii silných URL indikátorov, pričom má tendenciu označovať správy ako podozrivé aj v prípadoch, keď nie sú prítomné jednoznačné rizikové znaky.

Model strojového učenia vykazuje vyváženejšie správanie a jeho chybovosť nie je priamo viazaná na prítomnosť konkrétneho indikátora.

\paragraph{Zachovanie indikátorov}

Výsledky ukazujú, že URL indikátory zostávajú zachované aj po transporte správ cez experimentálne prostredie (GoPhish + MailHog). To potvrdzuje validitu experimentu a umožňuje spoľahlivé vyhodnotenie správania analyzátora v realistickom scenári.

\paragraph{Záver experimentu}

Experiment 3 potvrdil, že URL indikátory predstavujú významný zdroj informácií pre detekciu podozrivých správ. Heuristický analyzátor je veľmi citlivý na prítomnosť týchto znakov, avšak má tendenciu k vyššiemu počtu false positive.

Model strojového učenia poskytuje lepší celkový výkon a je vhodnejší pre finálne rozhodovanie. URL indikátory sú preto najvhodnejšie využiteľné ako doplnkový signál alebo vysvetľujúci faktor v rámci komplexnejšieho detekčného systému.

Výsledky Experimentu 3 zároveň potvrdzujú, že samotná prítomnosť URL indikátorov nie je postačujúca pre presné rozhodovanie a je potrebné ich kombinovať s komplexnejšími modelmi, ktoré dokážu lepšie zachytiť kontext správy.

\subsubsection{Experiment 4 – Robustnosť modelov voči distribučným zmenám a výber finálneho riešenia}\label{exp4}

Cieľom Experimentu 4 bolo systematicky vyhodnotiť robustnosť navrhnutých modelov
v rôznych scenároch distribučného posunu a na základe toho identifikovať
najvhodnejšiu konfiguráciu pre nasadenie v realistickom prostredí.
Experiment nadväzuje na predchádzajúce experimenty a rozširuje ich o
komplexné offline hodnotenie na viacerých dátových splitoch (Track A)
a validáciu na realistických kampaniach (Track B a Track C).

\paragraph{Track A – Offline robustnosť naprieč splitmi}

Modely boli vyhodnocované na piatich typoch splitov:
\texttt{val\_iid}, \texttt{test\_iid}, \texttt{test\_hard\_source},
\texttt{test\_hard\_cluster} a \texttt{test\_deployment}.
IID splity reprezentujú štandardné validačné nastavenie,
zatiaľ čo hard splity simulujú realistické distribučné posuny.

Výsledky ukazujú, že modely dosahujú stabilne vysoký výkon na IID dátach
(macro F1 $\approx$ 0.87--0.90), avšak výrazné rozdiely sa objavujú
na hard splitoch. V scenári L (\texttt{scenario\_l\_combined\_anti\_source})
dochádza k poklesu výkonu na \texttt{test\_hard\_source}
(macro F1 $\approx$ 0.69), čo naznačuje obmedzenú schopnosť
generalizácie na nové zdroje dát.

Naopak scenáre M vykazujú výrazne vyššiu robustnosť. V prípade rebalancovaného scenára
dosahuje model hodnoty macro F1 približne 0.84 na
\texttt{test\_hard\_source} a približne 0.90 na
\texttt{test\_deployment}, čo predstavuje významné zlepšenie
oproti scenáru L.

\begin{figure}[t]
\centering
\includegraphics[width=\linewidth]{exp4_trackA_deployment_macro_f1.pdf}
\caption{Porovnanie modelov na \texttt{test\_deployment}. Najlepší výkon
dosahuje scenár \texttt{m\_rebalanced\_fraud\_spam} s modelom SGD
(macro F1 = 0.901).}
\label{fig:exp4_trackA_deployment}
\end{figure}

Najlepší výsledok bol dosiahnutý v scenári
\texttt{m\_rebalanced\_fraud\_spam} s modelom
\texttt{hybrid\_sgd\_log}, ktorý dosahuje macro F1 = 0.901
na \texttt{test\_deployment}. Tento výsledok potvrdzuje, že kombinácia
anti-template prístupu a rebalancovania tried vedie k optimálnemu
kompromisu medzi presnosťou a robustnosťou.

\paragraph{Vplyv architektúry modelu}

Porovnanie modelov \texttt{hybrid\_logreg} a
\texttt{hybrid\_sgd\_log} ukazuje, že výkonnosť modelu závisí
od konkrétneho scenára. V anti-template nastavení dosahuje lepšie
výsledky logistická regresia, zatiaľ čo v rebalancovanom scenári
vykazuje lepšiu robustnosť model trénovaný pomocou SGD.
Tento výsledok naznačuje, že optimalizačný algoritmus interaguje
s charakterom dát a ovplyvňuje schopnosť modelu generalizovať.

\paragraph{Analýza tried}

Detailná analýza ukazuje, že trieda \texttt{phishing}
je detegovaná s veľmi vysokým recall naprieč všetkými scenármi,
čo naznačuje vysokú citlivosť modelu na typické phishingové vzory.
Naopak trieda \texttt{financial\_fraud} predstavuje najväčšiu výzvu,
pričom jej recall je v scenári L výrazne nižší, no v scenároch M
dochádza k jeho výraznému zlepšeniu (až nad 0.9).
Tento efekt je priamo spojený s rebalancovaním dát.

\paragraph{Track B – Doplnková validácia na realistickej kampani (Experiment 1)}

Ako hlavná realistická validácia modelu bola použitá kampaň
z Experimentu 1\ref{epx1}, ktorá najlepšie reprezentuje kompletný
end-to-end tok
(\textit{generovanie} $\rightarrow$ \textit{seeding} $\rightarrow$ \textit{zber}
$\rightarrow$ \textit{analýza} $\rightarrow$ \textit{klasifikácia})
na obsahovo zmiešaných dátach. Na rozdiel od špecializovaných kampaní
v Experimentoch 2 a 3 nejde o úzko zameraný benchmark na jednu doménu
príznakov (napr. iba hlavičky alebo iba URL), ale o realistickejší
vstup pre finálne ML rozhodovanie.

Výsledky z Track B potvrdzujú, že model vybraný na základe
offline robustnostného hodnotenia si zachováva vysokú kvalitu
aj v kampanovom prostredí. Táto časť preto slúži ako praktické
overenie prenositeľnosti offline výsledkov do realistického
nasadenia a dopĺňa štandardné hodnotenie na testovacích splitoch.

\paragraph{Track C – Multiclass doplnková analýza na špecializovaných kampaniach (Experimenty 2 a 3)}

Popri hlavnej kampanovej validácii v binárnom režime bol model
vyhodnotený aj v multiclass nastavení na kampaniach z Experimentu 2
(header-oriented) a Experimentu 3\ref{epx3} (URL-oriented).
Multiclass úloha zahŕňa triedy \texttt{legit}, \texttt{phishing},
\texttt{spam} a \texttt{financial\_fraud} a predstavuje doplnkovú
analytickú vrstvu, ktorá umožňuje detailnejšie posúdiť schopnosť
modelu rozlišovať medzi jednotlivými typmi útokov.

Na kampani z Experimentu~2 \ref{epx2} dosiahol model hodnoty
\begin{itemize}
    \item \texttt{mc\_accuracy = 0.51},
    \item \texttt{mc\_balanced\_accuracy = 0.5975},
    \item \texttt{mc\_f1\_macro = 0.5329}.
\end{itemize}

Na kampani z Experimentu~3 \ref{epx3} dosiahol
\begin{itemize}
    \item \texttt{mc\_accuracy = 0.52},
    \item \texttt{mc\_balanced\_accuracy = 0.6125},
    \item \texttt{mc\_f1\_macro = 0.5316}.
\end{itemize}


Pre kontext je vhodné uviesť vývoj multiclass výkonu v Experimente~1\ref{epx1},
kde sa rovnaká úloha postupne zlepšovala v rámci iterácií scenára~M.
Hodnota \texttt{mc\_f1\_macro}dosaihla až \texttt{0.7305}.
Tento trend potvrdzuje, že kombinácia textových reprezentácií
a ručne definovaných príznakov spolu s cieleným vyvážením tried
výrazne zlepšuje schopnosť modelu vykonávať jemnejšiu multiclass
diskrimináciu.

V porovnaní s binárnym hodnotením (F1 $>$ 0.96) je multiclass úloha
výrazne náročnejšia. Model spoľahlivo rozlišuje medzi legitímnymi
a podozrivými správami, avšak detailné rozlíšenie medzi jednotlivými
typmi útokov (najmä medzi \texttt{spam} a
\texttt{financial\_fraud}) zostáva limitované.

Pri interpretácii výsledkov je potrebné zohľadniť, že kampane z
Experimentov~2\ref{epx2} a~3\ref{epx3} sú zámerne konštruované so zameraním na konkrétne
typy príznakov (hlavičkové resp. URL indikátory). Nejde teda o
plnohodnotný všeobecný benchmark, ale o doplnkové testovanie
robustnosti modelu voči špecifickým distribučným posunom a typom útokov.

\paragraph{Výber finálneho modelu}

Na základe kombinácie výsledkov zo všetkých splitov a kampaní bol
ako finálny model zvolený
\texttt{scenario\_m\_rebalanced\_fraud\_spam} s modelom
\texttt{hybrid\_sgd\_log}. Tento model dosahuje najlepší kompromis
medzi presnosťou na \texttt{test\_deployment}, robustnosťou voči
distribučným posunom (\texttt{test\_hard\_source}) a schopnosťou
detegovať minoritné triedy.

\paragraph{Zhrnutie experimentu}

Experiment 4 ukázal, že hodnotenie modelov výlučne na IID dátach
nie je postačujúce pre realistické nasadenie. Scenáre zamerané
na variabilitu dát (anti-template) a vyváženie tried (rebalancing)
významne zlepšujú robustnosť modelov.

Výsledky zároveň naznačujú, že kombinácia vhodného tréningového
scenára a optimalizačnej metódy je kľúčová pre dosiahnutie
optimálneho výkonu. Navrhnutý finálny model predstavuje robustné
riešenie schopné generalizovať na nové typy útokov a praktické
nasadenie.

\subsubsection{Experiment 5 – Porovnanie heuristického, ML a hybridného prístupu na jednotnom benchmarku}\label{exp5}

Cieľom Experimentu 5 bolo systematicky porovnať tri prístupy k detekcii podozrivých e-mailov
– heuristický analyzátor, model strojového učenia a hybridné kombinácie – v jednotnom
a reprezentatívnom benchmarku. Experiment nadväzuje na predchádzajúce experimenty a
prepája ich do jedného hodnotiaceho rámca, ktorý umožňuje analyzovať robustnosť,
generalizáciu a správanie modelov naprieč rôznymi typmi dát.

\paragraph{Model setup}

Ako reprezentant modelového prístupu bol použitý model zo scenára
\texttt{scenario\_m\_rebalanced\_fraud\_spam} s klasifikátorom
\texttt{hybrid\_logreg} a vstupným režimom \texttt{text\_plus\_features}.
Tento model bol zvolený na základe výsledkov Experimentu 4 \ref{exp4}, kde preukázal
najlepší kompromis medzi výkonom a robustnosťou.

\paragraph{Zostavenie benchmarku}

Na rozdiel od predchádzajúcich experimentov bol vytvorený jednotný benchmark,
ktorý kombinuje tri typy dát:

\begin{itemize}
    \item \textbf{controlled benchmark} – kontrolované scenáre (Track A),
    \item \textbf{realistic campaign} – simulované phishing kampane (Track B),
    \item \textbf{offline generalization} – reálne distribuované dáta (\texttt{test\_deployment}).
\end{itemize}

Tento návrh umožňuje analyzovať nielen výkon modelu,
ale aj jeho stabilitu pri zmene distribúcie dát.

\paragraph{Porovnávané prístupy}

V experimente boli porovnané štyri prístupy:

\begin{itemize}
    \item \texttt{heuristic\_only} – pravidlový systém,
    \item \texttt{ml\_only} – model strojového učenia,
    \item \texttt{hybrid\_a\_conservative} – ML model doplnený heuristickým override,
    \item \texttt{hybrid\_b\_gated} – heuristika ako gating mechanizmus.
\end{itemize}

\paragraph{Celkové binárne výsledky}

\begin{table}[H]
\centering
\caption{Globálne porovnanie metód na jednotnom benchmarke}
\label{tab:global_comparison}
\footnotesize
\setlength{\tabcolsep}{4pt}
\renewcommand{\arraystretch}{1.1}
\resizebox{\columnwidth}{!}{%
\begin{tabular}{lcccc}
\toprule
\textbf{Metóda} & \textbf{Acc} & \textbf{Bal Acc} & \textbf{F1} & \textbf{MCC} \\
\midrule
\texttt{heuristic\_only}         & 0.76 & 0.73 & 0.69 & 0.51 \\
\texttt{ml\_only}                & \textbf{0.93} & \textbf{0.91} & \textbf{0.91} & \textbf{0.85} \\
\texttt{hybrid\_a\_conservative} & 0.92 & 0.90 & 0.90 & 0.83 \\
\texttt{hybrid\_b\_gated}        & 0.76 & 0.73 & 0.69 & 0.51 \\
\bottomrule
\end{tabular}%
}
\end{table}
Výsledky ukazujú, že model strojového učenia dosahuje najlepší celkový výkon
vo všetkých sledovaných metrikách. Hybrid A sa výkonom približuje ML modelu,
avšak neprináša výrazné zlepšenie.

Naopak hybrid B vykazuje prakticky identické výsledky ako heuristika,
čo je dôsledkom jeho návrhu – heuristika tu slúži ako rozhodovací filter,
a tým určuje finálne správanie systému.

\paragraph{Výsledky podľa typu dát}

\begin{table}[H]
\centering
\caption{Výkon metód podľa typu dát (F1 skóre)}
\label{tab:method_performance}
\footnotesize
\setlength{\tabcolsep}{4pt}
\renewcommand{\arraystretch}{1.1}
\resizebox{\columnwidth}{!}{%
\begin{tabular}{lcccc}
\toprule
\textbf{Skupina} & \textbf{Heuristika} & \textbf{ML} & \textbf{Hybrid A} & \textbf{Hybrid B} \\
\midrule
controlled benchmark   & 0.736 & 0.722 & \textbf{0.760} & 0.736 \\
realistic campaign     & 0.80  & \textbf{0.94}  & 0.87  & 0.80 \\
offline generalization & 0.04  & \textbf{0.93}  & 0.92  & 0.04 \\
\bottomrule
\end{tabular}%
}

\end{table}
Táto analýza odhaľuje zásadné rozdiely medzi prístupmi:

V kontrolovanom benchmarku dosahuje najlepší výkon hybrid A,
čo naznačuje, že heuristické pravidlá môžu byť prínosné
v prostredí s presne definovanými indikátormi.

V realistických kampaniach však heuristika stráca presnosť,
zatiaľ čo ML model dokáže efektívne generalizovať a zachytávať
komplexnejšie vzory.

Najvýraznejší rozdiel je v offline generalization,
kde heuristika prakticky zlyháva (F1 = 0.04),
čo potvrdzuje jej silnú závislosť na konkrétnych pravidlách
a nízku schopnosť adaptácie na nové dáta.

\paragraph{Multiclass výsledky}

Multiclass hodnotenie bolo realizované na podmnožine 1250 vzoriek,
pre ktoré bola dostupná ground truth klasifikácia.

\begin{table}[H]
\centering
\begin{tabular}{lc}
\hline
Metóda & Macro F1 \\
\hline
heuristic\_only & 0.333 \\
ml\_only & 0.679 \\
hybrid\_a\_conservative & 0.679 \\
hybrid\_b\_gated & 0.458 \\
\hline
\end{tabular}
\caption{Multiclass porovnanie metód (macro F1). Tabuľka ukazuje schopnosť modelov rozlišovať medzi typmi útokov, čo predstavuje náročnejší scenár než binárna klasifikácia.}
\end{table}

Heuristický prístup nedokáže efektívne rozlišovať medzi jednotlivými typmi útokov,
čo vedie k nízkemu výkonu. Model strojového učenia a hybrid A dosahujú výrazne
lepšie výsledky, keďže využívajú komplexnejšie reprezentácie dát.

Hybrid B vykazuje stredný výkon, čo opäť reflektuje jeho závislosť
na heuristickom gating mechanizme.

\paragraph{Interpretácia hybridných prístupov}

Hybrid A predstavuje konzervatívne rozšírenie ML modelu,
kde heuristika slúži ako dodatočný bezpečnostný mechanizmus.
Jeho prínos je však obmedzený a závisí od konkrétneho nastavenia prahov.

Hybrid B využíva heuristiku ako primárny filter,
čo vedie k tomu, že jeho správanie je takmer identické
s heuristickým prístupom. Tento výsledok ukazuje,
že nevhodne navrhnutá integrácia heuristík môže potlačiť
výhody modelu strojového učenia.

\paragraph{Záver experimentu}

Experiment 5 potvrdzuje, že model strojového učenia predstavuje
najrobustnejší prístup naprieč rôznymi typmi dát.

Heuristika je síce interpretovateľná a efektívna v kontrolovaných scenároch,
avšak výrazne zlyháva pri generalizácii.

Hybridné prístupy neprinášajú automatické zlepšenie výkonu,
a ich úspešnosť závisí od konkrétneho spôsobu integrácie.

Výsledky naznačujú, že heuristiky sú najvhodnejšie ako doplnkový
alebo vysvetľujúci komponent, zatiaľ čo finálne rozhodovanie
by malo byť založené na modeloch strojového učenia.


\subsubsection{Experiment 6 – Porovnanie nášho analyzátora s nástrojom PhishTool}

Cieľom Experimentu 6 bolo porovnať náš detekčný systém s externým nástrojom
PhishTool v realistickom analytickom scenári. Na rozdiel od predchádzajúcich
experimentov, ktoré sa zameriavali na automatizované vyhodnotenie modelov,
tento experiment reflektuje praktické použitie nástrojov pri manuálnej
analýze e-mailov.

\paragraph{Dataset a scenár}

Experiment bol realizovaný na menšom datasete (community-limit verzia),
ktorý obsahoval celkovo 25 e-mailov.

\begin{table}[!t]
\centering
\caption{Zloženie datasetu v Experimente 6}
\label{tab:exp6_dataset}
\footnotesize
\setlength{\tabcolsep}{4pt}
\renewcommand{\arraystretch}{1.1}
\begin{tabular}{lcc}
\toprule
\textbf{Kategória} & \textbf{Trieda} & \textbf{Počet} \\
\midrule
realistic campaign & legit & 7 \\
realistic campaign & phishing & 5 \\
realistic campaign & spam & 4 \\
realistic campaign & financial\_fraud & 4 \\
\midrule
controlled benchmark & spoofing\_header\_heavy & 2 \\
controlled benchmark & url\_heavy & 3 \\
\bottomrule
\end{tabular}
\end{table}

Takéto zloženie umožňuje kombinovať realistické scenáre s kontrolovanými
prípadmi obsahujúcimi špecifické technické indikátory.

Keďže nástroj PhishTool neposkytuje automatizované vyhodnotenie, ale slúži
primárne ako analytický a vizualizačný nástroj, bolo potrebné pri tvorbe
datasetu zvoliť vhodnú metodiku. Použitie plne kontrolovaného datasetu by
viedlo k skresleniu výsledkov, keďže by boli vopred známe použité indikátory
a typy manipulácií.

Z tohto dôvodu bol dataset konštruovaný dvojfázovo. Najprv bolo náhodne
vybraných 10 e-mailov z existujúcich datasetov (napr. Enron a interné datasety),
ktoré slúžili ako realistické referenčné vzorky. Následne boli tieto e-maily
použité ako vstupné príklady pre LLM, ktoré vygenerovalo ďalších 15 e-mailov.

Výsledný dataset tak obsahuje kombináciu reálnych a syntetických správ,
pričom pre syntetickú časť nie je explicitne známe presné rozdelenie tried
ani prítomné anomálie. Tento prístup umožňuje simulovať realistickejšie
podmienky analýzy, kde analytik nemá apriórnu znalosť o charaktere dát.

\paragraph{Workflow experimentu}

Experiment prebiehal v niekoľkých krokoch:

\begin{enumerate}
    \item výber reprezentatívnych \texttt{.eml} súborov,
    \item export do datasetu,
    \item item manuálna analýza pomocou nástroja PhishTool,
    \item automatická analýza pomocou nášho nástroja (IOC extrakcia),
    \item porovnanie výsledkov na úrovni:
    \begin{itemize}
        \item detegovaných indikátorov (IOC overlap),
        \item finálneho rozhodnutia (classification).
    \end{itemize}
\end{enumerate}

Manuálna analýza bola doplnená o kvalitatívne poznámky
čo umožnilo hlbšie pochopiť rozdiely medzi nástrojmi.

\paragraph{Decision-level porovnanie}

\begin{table}[H]
\centering
\caption{Porovnanie rozhodnutí medzi nástrojmi a ground truth}
\begin{tabular}{lcccc}
\hline
Porovnanie & Accuracy & Precision & Recall & F1 \\
\hline
Náš nástroj vs truth & 0.84 & 1.00 & 0.778 & 0.875 \\
PhishTool vs truth & 0.92 & 0.90 & 1.00 & 0.947 \\
\hline
\end{tabular}

\end{table}

Zároveň bola pozorovaná zhoda medzi nástrojmi na úrovni 0.76.

Výsledky ukazujú zásadný rozdiel v správaní oboch prístupov:

\begin{itemize}
    \item náš nástroj je \textbf{konzervatívnejší} – dosahuje vysokú precision (1.00),
    ale nižší recall,
    \item PhishTool je \textbf{citlivejší} – zachytí všetky útoky (recall = 1.00),
    ale generuje viac false positives.
\end{itemize}

Konkrétne prípady false positives v PhishTool boli pozorované
pri legitímnych e-mailoch (napr. exp6\_003, exp6\_008).

\paragraph{Analýza IOC indikátorov}

\begin{table}[H]
\centering
\caption{Zhoda detekcie URL a technických indikátorov}
\begin{tabular}{lc}
\hline
Indikátor & Zhoda \\
\hline
shortener\_url & 1.00 \\
ip\_url & 1.00 \\
obfuscated\_url & 1.00 \\
risky\_attachment & 1.00 \\
anchor\_mismatch & 0.80 \\
suspicious\_tld & 0.64 \\
display\_name\_spoof & 0.28 \\
\hline
\end{tabular}
\end{table}

Silná zhoda bola pozorovaná pri technických URL indikátoroch,
čo naznačuje, že tieto signály sú jednoznačné a dobre detegovateľné
oboma nástrojmi.

Naopak, najväčšie rozdiely sa objavili pri autentifikačných a spoofing indikátoroch.

\paragraph{Case studies}

Analýza konkrétnych prípadov ukázala, že najväčšie rozdiely
nastávajú pri komplexných útokoch kombinujúcich viac signálov,
najmä v kategóriách phishing a financial fraud.

Typické problematické prípady zahŕňali:

\begin{itemize}
    \item kombináciu legitímne vyzerajúcich a škodlivých URL,
    \item slabé alebo nejednoznačné autentifikačné signály,
    \item sociálne inžinierstvo bez jasného technického indikátora.
\end{itemize}

\paragraph{Interpretácia výsledkov}

Experiment ukazuje dva odlišné prístupy k detekcii:

\begin{itemize}
    \item náš analyzátor je presný a konzervatívny,
    minimalizuje false positives,
    \item naša analýza za pomoci PhishTool je agresívnejšia a citlivejšia,
    zachytáva viac útokov, ale za cenu vyššieho počtu chýb.
\end{itemize}

Z praktického hľadiska to znamená, že pre skúseného bezpečnostného
analytika predstavuje PhishTool silný nástroj, ktorý umožňuje detailnú
manuálnu analýzu jednotlivých indikátorov a hlavičiek e-mailu.
Poskytuje široké množstvo technických informácií, ktoré je možné
interpretovať v kontexte konkrétneho incidentu.

Na druhej strane, pre bežného používateľa alebo menej skúseného analytika
je takýto výstup náročnejší na interpretáciu, keďže vyžaduje dobré
porozumenie e-mailovým protokolom, autentifikačným mechanizmom
a typickým znakom phishingových útokov.

Navrhnutý analyzátor sa preto snaží tento problém riešiť tým, že
zjednodušuje výsledok do priamo interpretovateľnej podoby.
Používateľovi poskytuje okamžité vyhodnotenie spolu s vysvetlením,
bez potreby hlbšej expertnej analýzy jednotlivých indikátorov.

\paragraph{Záver experimentu}

Experiment 6 potvrdzuje, že oba prístupy majú svoje výhody.

PhishTool poskytuje detailnejšiu a citlivejšiu analýzu,
ktorá je vhodná pre bezpečnostných expertov.

Na druhej strane náš systém ponúka jednoduchšie,
konzistentné a automatizované riešenie,
ktoré je vhodnejšie pre bežných používateľov alebo
automatizované spracovanie e-mailov.

Najlepšie výsledky by bolo možné dosiahnuť kombináciou oboch prístupov,
kde automatizovaný systém slúži ako prvá línia detekcie
a manuálna analýza ako doplnkový krok pri zložitejších prípadoch.


\subsubsection{Finálna konfigurácia detekčného systému}

Na základe výsledkov experimentov 1 až 6 bola navrhnutá finálna konfigurácia
detekčného systému, ktorá predstavuje najlepší kompromis medzi presnosťou,
robustnosťou a praktickou použiteľnosťou.

\paragraph{Zvolená konfigurácia}

\begin{itemize}
    \item \textbf{Scenár dát:} \texttt{scenario\_m\_rebalanced\_fraud\_spam}
    \item \textbf{Model:} \texttt{hybrid\_logreg} (single-stage klasifikátor)
    \item \textbf{Vstupný režim:} \texttt{text\_plus\_features}
    \item \textbf{Rozhodovanie:}
    \begin{itemize}
        \item finálne rozhodnutie realizuje ML model,
        \item heuristika slúži ako doplnková vrstva (risk score, explainability).
    \end{itemize}
    \item \textbf{Výstup systému:}
    \begin{itemize}
        \item binárny: \texttt{legit} vs \texttt{suspicious} (operatívne použitie),
        \item multiclass: \texttt{legit/phishing/spam/financial\_fraud} (triage a prioritizácia).
    \end{itemize}
\end{itemize}

\paragraph{Zdôvodnenie výberu}

Výber finálnej konfigurácie vychádza z konzistentných výsledkov naprieč experimentami:

\begin{itemize}
    \item \textbf{Experiment 1 a 5:}
    kombinovaný vstup \texttt{text\_plus\_features} dosahuje stabilne lepší výkon
    než text-only alebo header-only prístup, čo potvrdzuje význam kombinácie
    obsahových a technických signálov.
    
    \item \textbf{Experiment 4:}
    scenáre typu M sú robustnejšie než scenáre L,
    najmä na náročných splitoch (\texttt{test\_hard} a \texttt{test\_deployment}),
    čo podporuje ich použitie v reálnom nasadení.
    
    \item \textbf{Experiment 5:}
    model strojového učenia (\texttt{ml\_only}) dosahuje najlepší globálny výkon,
    pričom hybridné prístupy neprinášajú automatické zlepšenie bez starostlivého návrhu.
    
    \item \textbf{Experiment 2 a 3:}
    heuristické indikátory sú veľmi presné na detekciu konkrétnych znakov,
    ale menej spoľahlivé ako finálny rozhodovací mechanizmus.
    
    \item \textbf{Experiment 6:}
    porovnanie s nástrojom PhishTool ukazuje,
    že explainability vrstva je dôležitá pre interpretáciu výsledkov,
    avšak finálne rozhodovanie by malo zostať na strane ML modelu.
\end{itemize}

\paragraph{Zhrnutie}

Výsledná architektúra systému reflektuje kombinovaný prístup,
kde model strojového učenia predstavuje hlavný rozhodovací komponent
a heuristika slúži ako doplnkový zdroj informácií.

Takýto návrh umožňuje dosiahnuť vysokú presnosť pri zachovaní
interpretovateľnosti výsledkov, čo je kľúčové pre praktické nasadenie.

\subsubsection*{Limitácie a implikácie pre prax}

Napriek dosiahnutým výsledkom má navrhnutý systém viacero limitácií,
ktoré je potrebné zohľadniť pri jeho interpretácii a nasadení.

\begin{itemize}
    \item \textbf{Prekryv tried v multiclass úlohe:}
    triedy \texttt{phishing}, \texttt{spam} a \texttt{financial\_fraud}
    vykazujú výrazný štýlový prekryv,
    čo vedie k chybám pri hraničných prípadoch a znižuje presnosť kategorizácie.

    \item \textbf{Závislosť hlavičkových signálov od transportu:}
    v experimentálnom prostredí (GoPhish/MailHog) dochádza k normalizácii
    alebo strate autentifikačných hlavičiek,
    čo znamená, že výsledky z kontrolovaných benchmarkov nie sú
    plne prenositeľné do reálnej prevádzky.

    \item \textbf{Obmedzenia heuristickej vrstvy:}
    heuristika je citlivá na prítomnosť rizikových indikátorov,
    avšak pri URL-heavy alebo urgentných správach generuje zvýšený počet
    false positives, čo obmedzuje jej použitie ako samostatného rozhodovacieho nástroja.

    \item \textbf{Limitovaný rozsah externého porovnania:}
    Experiment 6 bol realizovaný na menšej vzorke dát
    a zahŕňal manuálnu anotáciu,
    čo obmedzuje štatistickú významnosť výsledkov.

    \item \textbf{Citlivosť na zmenu distribúcie dát (concept drift):}
    model je závislý od aktuálnych jazykových vzorov a techník útočníkov,
    preto vyžaduje pravidelný re-tréning a aktualizáciu.

\end{itemize}

\paragraph{Implikácie pre prax}

Z praktického pohľadu možno navrhnutý systém interpretovať ako
kombináciu viacerých komplementárnych komponentov. Jadrom riešenia
je robustný rozhodovací mechanizmus založený na strojovom učení,
ktorý zabezpečuje automatizovanú detekciu podozrivých správ.
Tento mechanizmus je doplnený heuristickou vrstvou, ktorá poskytuje
interpretovateľné vysvetlenie rozhodnutí a umožňuje lepšie pochopenie
identifikovaných rizík. Celkovo tak systém funguje ako podporný nástroj,
ktorý je použiteľný nielen pre bezpečnostných analytikov, ale aj pre
bežných používateľov bez hlbších odborných znalostí.

Získané výsledky naznačujú, že pre bežného používateľa je vhodnejšie
využiť automatizovaný analyzátor, ktorý poskytuje okamžité,
konzistentné a jednoducho interpretovateľné rozhodnutia bez potreby
detailnej znalosti problematiky. Naopak, pre skúseného bezpečnostného
analytika predstavujú nástroje ako PhishTool hodnotný doplnok, ktorý
umožňuje detailnú manuálnu analýzu jednotlivých indikátorov a hlbšie
preskúmanie konkrétnych prípadov.

Z pohľadu praktického nasadenia sa ako najefektívnejší javí kombinovaný
prístup, v ktorom automatizovaný systém slúži ako prvá línia detekcie
a filtrácie potenciálne rizikových správ, zatiaľ čo manuálna analýza
je využívaná ako následný verifikačný a investigatívny krok.

\paragraph{Záverečné zhodnotenie}

Experimenty 1 až 6 ukazujú, že kombinácia modelov strojového učenia
a heuristických indikátorov predstavuje efektívny prístup k detekcii
phishingových a podvodných e-mailov.

Kľúčovým poznatkom je, že heuristika by nemala byť použitá ako hlavný
rozhodovací mechanizmus, ale ako doplnkový a vysvetľujúci komponent,
ktorý zvyšuje dôveru v rozhodnutia modelu.

Takto navrhnutý systém je schopný dosiahnuť vysokú presnosť,
zachovať interpretovateľnosť a zároveň zostať prakticky použiteľný
v reálnom prostredí.


\section{Záver}

Táto práca sa zaoberala problematikou bezpečnosti e-mailovej komunikácie, analýzou najčastejších hrozieb a návrhom systému na detekciu podozrivých e-mailov. V teoretickej časti boli predstavené princípy fungovania e-mailovej komunikácie, formát správ a techniky využívané pri útokoch, ako sú phishing, spam alebo spoofing, spolu s mechanizmami ochrany.

V praktickej časti bol navrhnutý a implementovaný systém kombinujúci heuristický analyzátor a prístupy strojového učenia. Heuristický analyzátor umožňuje transparentnú a interpretovateľnú analýzu e-mailov na základe technických a obsahových indikátorov, zatiaľ čo modely strojového učenia poskytujú vyššiu presnosť a schopnosť generalizácie.

Experimenty ukázali, že samotné heuristické prístupy dosahujú vysoký recall, avšak modelové prístupy poskytujú lepší kompromis medzi presnosťou a úplnosťou detekcie. Najlepšie výsledky boli dosiahnuté kombináciou textovej analýzy a odvodených príznakov, pričom finálny model preukázal dobrú schopnosť rozlišovať medzi legitímnymi a škodlivými e-mailmi aj v podmienkach blízkych reálnemu nasadeniu.

Významným prínosom práce je návrh experimentálneho rámca s dôrazom na elimináciu problémov ako source leakage a template leakage, čo umožnilo realistickejšie hodnotenie modelov a ich robustnosti. Tento prístup zvyšuje dôveryhodnosť dosiahnutých výsledkov a ich prenositeľnosť do praxe.

Do budúcnosti je možné prácu rozšíriť o využitie pokročilejších modelov hlbokého učenia, integráciu systému do reálneho prostredia e-mailových služieb a rozšírenie datasetov o nové typy útokov. Zároveň by bolo možné skúmať adaptívne a online učenie modelov na dynamicky sa meniace hrozby.

\begin{thebibliography}{00}

\bibitem{rfc5598}
D. Crocker, ``Internet Mail Architecture,'' RFC 5598, Internet Engineering Task Force (IETF), Jul. 2009. [Online]. Available: https://www.rfc-editor.org/rfc/rfc5598.html

\bibitem{rfc5321}
J. Klensin, ``Simple Mail Transfer Protocol,'' RFC 5321, Internet Engineering Task Force (IETF), Oct. 2008. [Online]. Available: https://datatracker.ietf.org/doc/html/rfc5321

\bibitem{rfc5322}
P. Resnick, ``Internet Message Format,'' RFC 5322, Internet Engineering Task Force (IETF), Oct. 2008. [Online]. Available: https://datatracker.ietf.org/doc/html/rfc5322

\bibitem{rfc2045}
N. Freed and N. Borenstein, ``Multipurpose Internet Mail Extensions (MIME) Part One: Format of Internet Message Bodies,'' RFC 2045, Internet Engineering Task Force (IETF), Nov. 1996. [Online]. Available: https://www.rfc-editor.org/rfc/rfc2045

\bibitem{rfc2046}
N. Freed and N. Borenstein, ``Multipurpose Internet Mail Extensions (MIME) Part Two: Media Types,'' RFC 2046, Internet Engineering Task Force (IETF), Nov. 1996. [Online]. Available: https://datatracker.ietf.org/doc/html/rfc2046

\bibitem{karim2019spam}
A. Karim, S. Azam, B. Shanmugam, K. Kannoorpatti, and M. Jonkman, ``A comprehensive survey for intelligent spam email detection,'' IEEE Access, vol. 7, pp. 168261–168295, 2019. [Online]. Available: https://ieeexplore.ieee.org/document/8907831

\bibitem{jakobsson2007phishing}
M. Jakobsson and S. Myers, \textit{Phishing and Countermeasures: Understanding the Increasing Problem of Electronic Identity Theft}. Wiley-Interscience, 2006. [Online]. Available: https://download.e-bookshelf.de/download/0000/5686/57/L-G-0000568657-0002356793.pdf

\bibitem{aurangzeb2017ransomware}
K. Aurangzeb, M. Aleem, M. Iqbal, and M. Islam, ``Ransomware: A Survey and Trends,'' 2017. [Online]. Available: https://www.researchgate.net/publication/317380115\_Ransomware\_A\_Survey\_and\_Trends

\bibitem{verma2017phishing}
J. Stachowski, ``Analysis of Phishing Emails and Detection Techniques,'' M.S. thesis, St. Cloud State University, 2017. [Online]. Available: https://repository.stcloudstate.edu/msia\_etds/125/
\bibitem{rfc7208}
S. Kitterman, ``Sender Policy Framework (SPF) for Authorizing Use of Domains in Email, Version 1,'' RFC 7208, Internet Engineering Task Force (IETF), Apr. 2014. [Online]. Available: https://datatracker.ietf.org/doc/html/rfc7208

\bibitem{rfc6376}
M. Crocker, D. Hansen, and T. Kucherawy, ``DomainKeys Identified Mail (DKIM) Signatures,'' RFC 6376, Internet Engineering Task Force (IETF), Sep. 2011. [Online]. Available: https://datatracker.ietf.org/doc/html/rfc6376

\bibitem{mailhog}
I. Charlie, ``MailHog: Email testing tool for developers,'' GitHub repository. [Online]. Available: https://github.com/mailhog/MailHog

\bibitem{gophish}
J. Wright, ``GoPhish: Open-source phishing framework,'' GitHub repository. [Online]. Available: https://github.com/gophish/gophish

\bibitem{phishtool}
PhishTool, ``PhishTool: Platform for phishing analysis and response,'' [Online]. Available: https://www.phishtool.com

\bibitem{chatgpt}
OpenAI, ``ChatGPT,'' 2026. [Online]. Available: https://chat.openai.com. 
Použité ako podporný nástroj pri úprave textu (gramatika, štylistika) a asistencii pri tvorbe programového riešenia.

\bibitem{regexurl}
J. D'Alessandro, ``Comprehensive Regex for URL Detection and Spam Filtering,'' 2023. [Online]. Available: https://johndalesandro.com/blog/comprehensive-regex-for-url-detection-and-spam-filtering/. 
Použité ako inšpirácia pri návrhu regulárnych výrazov na detekciu URL adries v e-mailových správach.
\end{thebibliography}

\vspace{12pt}


\end{document}
